% =========================================================================
% SciPost LaTeX template
% Version 2024-07
%
% Submissions to SciPost Journals should make use of this template.
%
% INSTRUCTIONS: simply look for the `TODO:' tokens and adapt your file.
% ========================================================================

\documentclass{SciPost}

% Prevent all line breaks in inline equations.
\binoppenalty=10000
\relpenalty=10000

\hypersetup{
    colorlinks,
    linkcolor={red!50!black},
    citecolor={blue!50!black},
    urlcolor={blue!80!black}
}

\usepackage[bitstream-charter]{mathdesign}
\urlstyle{same}

% Number lines in verbatim code blocks for easy reference
\usepackage{fancyvrb}
\RecustomVerbatimEnvironment{verbatim}{Verbatim}{numbers=left,numbersep=5pt,frame=single,fontsize=\small}

% Fix \cal and \mathcal characters look (so it's not the same as \mathscr)
\DeclareSymbolFont{usualmathcal}{OMS}{cmsy}{m}{n}
\DeclareSymbolFontAlphabet{\mathcal}{usualmathcal}

\fancypagestyle{SPstyle}{
\fancyhf{}
\lhead{\colorbox{scipostblue}{\bf \color{white} ~SciPost Physics Core }}
\rhead{{\bf \color{scipostdeepblue} ~Submission }}
\renewcommand{\headrulewidth}{1pt}
\fancyfoot[C]{\textbf{\thepage}}
}

\begin{document}

\pagestyle{SPstyle}

\begin{center}{\Large \textbf{\color{scipostdeepblue}{
%%%%%%%%%% TODO: Write your article's title here
SptSet GAP Library Introduction
%%%%%%%%%% END TODO: TITLE
}}}\end{center}

\begin{center}\textbf{
%%%%%%%%%% TODO: AUTHORS
% Write the author list here. 
% Use (full) first name (+ middle name initials) + surname format.
% Separate subsequent authors by a comma, omit comma and use "and" for the last author.
% Mark the corresponding author(s) with a superscript symbol in this order
% \star, \dagger, \ddagger, \circ, \S, \P, \parallel, ...
Xing-Yu Ren\textsuperscript{1$\star$}
%%%%%%%%%% END TODO: AUTHORS
}\end{center}

\begin{center}
%%%%%%%%%% TODO: AFFILIATIONS
% Write all affiliations here.
% Format: institute, city, country
{\bf 1} Department of Physics, The Chinese University of Hong Kong
%%%%%%%%%% END TODO: AFFILIATIONS
%%%%%%%%%% TODO: EMAIL
% Provide email address of corresponding author(s)
\\ [\baselineskip]
$\star$ \href{mailto:xy.ren@link.cuhk.edu.hk}{\small xy.ren@link.cuhk.edu.hk}\
%%%%%%%%%% END TODO: EMAIL
\end{center}


\section*{\color{scipostdeepblue}{Abstract}}
\textbf{\boldmath{The goal of this note is to give a very detailed manual for the SptSet package built on GAP system.}}



%%%%%%%%%% TODO: LINENO
% For convenience during refereeing we turn on line numbers:
\linenumbers
% You should run LaTeX twice in order for the line numbers to appear.
%%%%%%%%%% END TODO: LINENO

%%%%%%%%%% TODO: TOC 
% Guideline: if your paper is longer that 6 pages, include a TOC
% To remove the TOC, simply cut the following block
\vspace{10pt}
\noindent\rule{\textwidth}{1pt}
\tableofcontents
\noindent\rule{\textwidth}{1pt}
\vspace{10pt}
%%%%%%%%%% END TODO: TOC


%%%%%%%%% TODO: CONTENTS 

\section{Overview: A Quick Practical Introduction}
\label{sec:overview}

This section provides a rapid introduction to the SptSet package for readers who understand the physics of SPT/SET phases but are new to GAP programming and this codebase.

\subsection{What is SptSet?}

SptSet is a GAP package for computing the classification of Symmetry Protected Topological (SPT) and Symmetry Enriched Topological (SET) phases. The package implements spectral sequence methods to compute group cohomology with twisted coefficients, which classifies these phases.

\textbf{Current capabilities:}
\begin{itemize}
    \item 2D fermion SPT classification (fully working)
    \item 3D classification when symmetry is $G_b \times \mathbb{Z}_2^f$ (direct product)
    \item Support for spin-1/2 systems
\end{itemize}

\subsection{GAP Basics for SptSet Users}

Before using SptSet, you need minimal GAP knowledge:

\subsubsection{GAP Syntax Essentials}
\begin{verbatim}
# Comments start with #
x := 5;                    # Variable assignment
[1, 2, 3]                  # A list
rec(a := 1, b := 2)        # A record (like dictionary)
f := x -> x^2;             # Lambda function
f := function(x) return x^2; end;  # Full function syntax
\end{verbatim}

\subsubsection{Package File Conventions}
\begin{itemize}
    \item \texttt{*.gd} files: \textbf{Declarations} -- define what functions exist
    \item \texttt{*.gi} files: \textbf{Implementations} -- define how functions work
    \item \texttt{init.g}: Loads all \texttt{.gd} files when package loads
    \item \texttt{read.g}: Loads all \texttt{.gi} files when package loads
\end{itemize}

\subsection{Package Architecture}

The SptSet package has a layered architecture:

\begin{center}
\begin{tabular}{|c|l|}
\hline
\textbf{Layer} & \textbf{Modules} \\
\hline
User Interface & \texttt{ss\_fermion.g*}, \texttt{ss\_fermion\_ez.g*}, etc. \\
\hline
Spectral Sequence & \texttt{spectral\_sequence.g*}, \texttt{ss\_vanilla.g*} \\
\hline
Cohomology & \texttt{group\_cohomology.g*}, \texttt{cochain.g*} \\
\hline
Coefficients & \texttt{coefficient.g*}, \texttt{module.g*} \\
\hline
Resolution & \texttt{bar\_resolution\_map*.g*} (interfaces with HAP) \\
\hline
\end{tabular}
\end{center}

\subsection{Running Your First Example}
\label{subsec:first-example}

This subsection provides a complete walkthrough of running an SPT classification, with line-by-line explanations of both the input commands and their outputs.

\subsubsection{Physical Background}

In fermionic systems, the total symmetry group $G_f$ is a $\mathbb{Z}_2$ extension of the bosonic symmetry $G_b$:
\begin{equation}
1 \to \mathbb{Z}_2^f \to G_f \to G_b \to 1
\end{equation}
where $\mathbb{Z}_2^f = \{1, P_f\}$ is generated by the fermion parity operator $P_f = (-1)^F$. The fermion parity map tells us which elements of $G_b$ are ``fermionic'' (they anticommute with fermion operators when lifted to $G_f$).

For crystallographic space groups, there is a natural fermion parity assignment: point group operations with determinant $-1$ (reflections, improper rotations) are fermionic, while those with determinant $+1$ (rotations, translations) are bosonic.

\subsubsection{Step 1: Start GAP and Load the Package}
\begin{verbatim}
$ cd /opt/gap-4.15.1
$ ./gap
gap> LoadPackage("SptSet");
\end{verbatim}

\subsubsection{Step 2: Set Up a Symmetry Group}

\paragraph{Input:}
\begin{verbatim}
gap> SG := SpaceGroupBBNWZ(2, 2);
\end{verbatim}

\paragraph{Output:}
\begin{verbatim}
SpaceGroupOnRightBBNWZ( 2, 1, 2, 1, 1 )
\end{verbatim}

\paragraph{Explanation:}
\begin{itemize}
    \item \texttt{SpaceGroupBBNWZ(d, n)}: A function from the \texttt{cryst} package that returns the $n$-th crystallographic space group in $d$ dimensions, using the BBNWZ (Brown-B\"ulow-Neub\"user-Wondratschek-Zassenhaus) classification.
    \item \texttt{2} (first argument): Dimension. For 2D, these are the 17 \textit{wallpaper groups}.
    \item \texttt{2} (second argument): Group number in the BBNWZ list. Valid range is 1--17 for 2D.
    \item The output \texttt{SpaceGroupOnRightBBNWZ( 2, 1, 2, 1, 1 )} is the internal representation. The parameters encode: dimension, crystal system, point group number, and additional identifiers for the space group type.
    \item \texttt{SG}: The result is an \textit{affine crystallographic group} object, where each element is represented as an affine transformation $x \mapsto Mx + t$ with $M$ an integer matrix (the point group part) and $t$ a translation vector.
\end{itemize}

\paragraph{Input:}
\begin{verbatim}
gap> fSG := IsomorphismPcpGroup(SG);
\end{verbatim}

\paragraph{Output:}
\begin{verbatim}
[ [ [ -1, 0, 0 ], [ 0, -1, 0 ], [ 0, 0, 1 ] ], 
  [ [ 1, 0, 0 ], [ 0, 1, 0 ], [ 1, 0, 1 ] ], 
  [ [ 1, 0, 0 ], [ 0, 1, 0 ], [ 0, 1, 1 ] ] ] -> [ g1, g2, g3 ]
\end{verbatim}

\paragraph{Explanation:}
\begin{itemize}
    \item \texttt{IsomorphismPcpGroup}: Computes an isomorphism from the input group to a \textit{polycyclic presentation} (Pcp) group.
    \item \textbf{Why Pcp?} The HAP package (which computes group cohomology resolutions) requires groups in Pcp format because it uses the polycyclic structure to efficiently compute resolutions.
    \item The output shows the mapping of three generators:
    \begin{itemize}
        \item Generator 1: $\begin{pmatrix} -1 & 0 & 0 \\ 0 & -1 & 0 \\ 0 & 0 & 1 \end{pmatrix} \mapsto \texttt{g1}$. This is a $3\times 3$ augmented matrix representing the affine transformation. The upper-left $2\times 2$ block is $\begin{pmatrix} -1 & 0 \\ 0 & -1 \end{pmatrix}$ (point inversion), and the third column $\begin{pmatrix} 0 \\ 0 \end{pmatrix}$ is the translation. The determinant of the point group part is $(-1)(-1) = +1$.
        \item Generator 2: Translation by $(1, 0)$. The identity matrix with translation vector. Determinant $= +1$.
        \item Generator 3: Translation by $(0, 1)$. Determinant $= +1$.
    \end{itemize}
    \item \texttt{fSG}: This is a \textit{group homomorphism} object. You can apply it to elements of \texttt{SG} using \texttt{Image(fSG, g)} to get the corresponding element in the Pcp group.
\end{itemize}

\paragraph{Input:}
\begin{verbatim}
gap> SG1 := Image(fSG);
\end{verbatim}

\paragraph{Output:}
\begin{verbatim}
Pcp-group with orders [ 2, 0, 0 ]
\end{verbatim}

\paragraph{Explanation:}
\begin{itemize}
    \item \texttt{Image(fSG)}: When called with just the homomorphism (no second argument), returns the \textit{image} of the homomorphism, i.e., the codomain group.
    \item \texttt{Pcp-group with orders [ 2, 0, 0 ]}: This describes the polycyclic structure:
    \begin{itemize}
        \item \texttt{g1} has order 2 (it squares to the identity---point inversion applied twice)
        \item \texttt{g2} has order 0 (infinite order---translation)
        \item \texttt{g3} has order 0 (infinite order---translation)
    \end{itemize}
    \item Mathematically, $G \cong \mathbb{Z}_2 \ltimes \mathbb{Z}^2$, where $\mathbb{Z}_2$ is the point group and $\mathbb{Z}^2$ is the translation lattice.
\end{itemize}

\subsubsection{Step 3: Compute a Resolution}

\paragraph{Input:}
\begin{verbatim}
gap> R := ResolutionAlmostCrystalGroup(SG1, 6);
\end{verbatim}

\paragraph{Output:}
\begin{verbatim}
Resolution of length 6 in characteristic 0 for 
Pcp-group with orders [ 2, 0, 0 ] .
\end{verbatim}

\paragraph{Explanation:}
\begin{itemize}
    \item \texttt{ResolutionAlmostCrystalGroup(G, n)}: A function from the HAP package that computes a free $\mathbb{Z}G$-resolution of $\mathbb{Z}$ up to degree $n$.
    \item \texttt{6}: The maximum dimension for the resolution. For 2D SPT classification, we need at least degree 3; for 3D, at least degree 4. Using 6 provides headroom.
    \item \texttt{characteristic 0}: The resolution is over $\mathbb{Z}$ (not a finite field).
    \item The resolution $R$ is a chain complex:
    \begin{equation}
    \cdots \to R_n \xrightarrow{d_n} R_{n-1} \to \cdots \to R_1 \xrightarrow{d_1} R_0 \xrightarrow{\epsilon} \mathbb{Z} \to 0
    \end{equation}
    where each $R_n$ is a free $\mathbb{Z}G$-module. This is the key algebraic input for computing group cohomology $H^n(G, M)$.
\end{itemize}

\subsubsection{Step 4: Define Fermion Parity}

\paragraph{Input:}
\begin{verbatim}
gap> gs := GeneratorsOfGroup(SG);
\end{verbatim}

\paragraph{Output:}
\begin{verbatim}
[ [ [ -1, 0, 0 ], [ 0, -1, 0 ], [ 0, 0, 1 ] ], 
  [ [ 1, 0, 0 ], [ 0, 1, 0 ], [ 1, 0, 1 ] ], 
  [ [ 1, 0, 0 ], [ 0, 1, 0 ], [ 0, 1, 1 ] ] ]
\end{verbatim}

\paragraph{Explanation:}
\begin{itemize}
    \item \texttt{GeneratorsOfGroup(SG)}: Returns the generators in matrix form.
    \item Each generator $g$ is a $(d+1) \times (d+1)$ augmented matrix encoding the affine map:
    \begin{equation}
    g = \begin{pmatrix} M & t \\ 0 & 1 \end{pmatrix}
    \end{equation}
    where $M$ is the $d \times d$ point group part and $t$ is the translation.
    \item \textbf{Important}: We use generators of \texttt{SG} (not \texttt{SG1}) because we need the matrix representation to compute determinants. Using \texttt{GeneratorsOfGroup(SG1)} would give Pcp elements, which cannot be passed to \texttt{DeterminantMat}.
\end{itemize}

\paragraph{Input:}
\begin{verbatim}
gap> f := GroupHomomorphismByImagesNC(SG1, GL(1, Integers),
>   List(gs, x -> Image(fSG, x)),
>   List(gs, x -> [[DeterminantMat(x)]]));
\end{verbatim}

\paragraph{Output:}
\begin{verbatim}
[ g1, g2, g3 ] -> [ [ [ 1 ] ], [ [ 1 ] ], [ [ 1 ] ] ]
\end{verbatim}

\paragraph{Explanation:}
This constructs the fermion parity homomorphism $f: G_b \to \mathbb{Z}_2 \cong \{+1, -1\}$.

\textbf{Function structure:} \texttt{GroupHomomorphismByImagesNC(Source, Target, SourceGens, TargetImages)}

\begin{itemize}
    \item \texttt{GroupHomomorphismByImagesNC}: Creates a group homomorphism by specifying where generators map to. The ``NC'' stands for ``No Check''---it skips verification that the map is well-defined.
    
    \item \texttt{SG1}: The \textbf{source group} (domain). Our space group in Pcp format.
    
    \item \texttt{GL(1, Integers)}: The \textbf{target group} (codomain). The group of $1 \times 1$ invertible integer matrices, $\{[[1]], [[-1]]\} \cong \mathbb{Z}_2$.
    
    \item \texttt{List(gs, x -> Image(fSG, x))}: The \textbf{source generators} in \texttt{SG1}.
    \begin{itemize}
        \item \texttt{List(collection, func)}: GAP's map function---applies \texttt{func} to each element.
        \item \texttt{x -> Image(fSG, x)}: A lambda function converting each matrix generator to its Pcp form.
        \item Result: \texttt{[ g1, g2, g3 ]}.
    \end{itemize}
    
    \item \texttt{List(gs, x -> [[DeterminantMat(x)]])}: The \textbf{target images}.
    \begin{itemize}
        \item \texttt{DeterminantMat(x)}: Computes determinant of the full augmented matrix. For space groups, this equals the determinant of the point group part (since the augmented matrix has 1 in the bottom-right corner).
        \item \texttt{[[...]]}: Double brackets wrap the scalar as a $1 \times 1$ matrix for \texttt{GL(1, Integers)}.
    \end{itemize}
\end{itemize}

\textbf{Output interpretation:} All three generators map to $[[1]]$, meaning all have determinant $+1$:
\begin{itemize}
    \item \texttt{g1}: Point inversion has $\det = (-1)(-1) = +1$
    \item \texttt{g2}, \texttt{g3}: Translations have $\det = 1$
\end{itemize}
This means the fermion parity is \textit{trivial} for this space group---there are no orientation-reversing elements. In physical terms, we are classifying SPT phases where all symmetry operations are bosonic.

\subsubsection{Step 5: Run the Classification}

\paragraph{Input:}
\begin{verbatim}
gap> SS := FermionEZSPTSpecSeq(R, f);
\end{verbatim}

\paragraph{Output:}
\begin{verbatim}
<object>
\end{verbatim}

\paragraph{Explanation:}
\begin{itemize}
    \item \texttt{FermionEZSPTSpecSeq(R, f)}: Creates a spectral sequence object for fermionic SPT classification.
    \item ``EZ'' stands for ``easy''---this is the simplified version where fermion parity is trivial or simple.
    \item The spectral sequence is a computational tool that organizes the cohomology calculation into ``layers'' that converge to the final classification.
    \item \texttt{<object>}: GAP's default display for complex objects. The actual data is stored internally.
\end{itemize}

\paragraph{Input:}
\begin{verbatim}
gap> FermionSPTLayersVerbose(SS, 2);
\end{verbatim}

\paragraph{Output:}
\begin{verbatim}
Majorana:
<ZL-Module with torsions [ 2, 2, 2 ]>
<ZL-Module with torsions [ 2, 2, 2 ]>
<ZL-Module with torsions [ 2, 2, 2 ]>
Complex fermion:
<ZL-Module with torsions [ 2, 2, 2, 2 ]>
<ZL-Module with torsions [ 2, 2, 2, 2 ]>
Bosonic:
<ZL-Module with torsions [ 2, 2, 2, 2 ]>
<ZL-Module with torsions [ 2, 2, 2, 2 ]>
<ZL-Module with torsions [ 2, 2, 2, 2 ]>
\end{verbatim}

\paragraph{Explanation:}
This is the main classification result. The function computes SPT phases up to dimension 2, organized by the type of fermionic degrees of freedom.

\textbf{Structure of the output:}
\begin{itemize}
    \item \textbf{Majorana}: The ``Majorana'' heading corresponds to the $q=2$ sector of the fermionic spectral sequence.
    
    	\textbf{Important:} the three displayed lines are \emph{not} three different physical ``layers''; they are successive \emph{spectral sequence pages} $E_r^{p,q}$ at fixed bidegree $(p,q)$ as $r$ increases.
    \begin{itemize}
        \item \texttt{<ZL-Module with torsions [ 2, 2, 2 ]>}: This is $\mathbb{Z}_2 \times \mathbb{Z}_2 \times \mathbb{Z}_2 = \mathbb{Z}_2^3$.
        \item In \texttt{FermionSPTLayersVerbose(SS, 2)} the code fixes total degree $p+q=2+1=3$ and prints the components $E_r^{p,q}$ for the relevant pairs $(p,q)$ with that total degree.
        \item Here ``Majorana'' corresponds to $(p,q)=(1,2)$, and the code prints $E_2^{1,2}$, $E_3^{1,2}$, and $E_4^{1,2}$ (three lines) to show how this component evolves under the differentials.
        \item The fact that all three lines are identical means this component stabilizes early (the relevant differentials vanish in this example).
    \end{itemize}
    
    \item \textbf{Complex fermion}: The $q=1$ sector. For \texttt{dim=2}, this corresponds to $(p,q)=(2,1)$, and the code prints $E_2^{2,1}$ and $E_3^{2,1}$ (two lines).
    \begin{itemize}
        \item \texttt{<ZL-Module with torsions [ 2, 2, 2, 2 ]>}: This is $\mathbb{Z}_2^4$.
        \item Again, the two lines are two pages ($r=2,3$), not two separate degrees.
    \end{itemize}
    
    \item \textbf{Bosonic}: The $q=0$ sector. For \texttt{dim=2}, this corresponds to $(p,q)=(3,0)$ and the code prints $E_2^{3,0}$, $E_3^{3,0}$, and $E_4^{3,0}$.
    \begin{itemize}
        \item \texttt{<ZL-Module with torsions [ 2, 2, 2, 2 ]>}: $\mathbb{Z}_2^4$ per page shown here.
        \item The three lines are the pages $r=2,3,4$ at fixed $(p,q)$.
    \end{itemize}
\end{itemize}

\textbf{Physical interpretation:}
\begin{itemize}
    \item Each \texttt{<ZL-Module with torsions [...] >} is the abelian group of a specific spectral-sequence component $E_r^{p,q}$ at a specific page $r$.
    \item Torsions like \texttt{[ 2, 2, 2 ]} mean the group is $\mathbb{Z}_2 \times \mathbb{Z}_2 \times \mathbb{Z}_2$.
    \item The spectral sequence provides a \emph{filtration} of the final classification group. Roughly, the stabilized groups (informally ``$E_\infty$'') give the associated-graded pieces, and the true classification is built from them via extension data.
    \item \textbf{Rule of thumb:} if the page-to-page prints stop changing, that component has stabilized (at least up to the printed cutoff), so you can treat that printed group as the final associated-graded piece for that $(p,q)$.
    \item For the wallpaper group \#2 with trivial fermion parity, we see rich $\mathbb{Z}_2$ torsion structure reflecting the topology of the quotient space and the group cohomology.
\end{itemize}

\subsection{Key Functions Reference}

\begin{center}
\begin{tabular}{|l|p{8cm}|}
\hline
\textbf{Function} & \textbf{Purpose} \\
\hline
\texttt{LoadPackage("SptSet")} & Load the SptSet package \\
\hline
\texttt{SpaceGroupBBNWZ(d, n)} & Get $d$-dimensional space group number $n$ \\
\hline
\texttt{ResolutionAlmostCrystalGroup(G, n)} & Compute resolution up to degree $n$ \\
\hline
\texttt{FermionEZSPTSpecSeq(R, f)} & Create spectral sequence for ``easy'' fermion parity \\
\hline
\texttt{FermionSPTSpecSeq(R, f, w)} & Create spectral sequence with spin structure $w$ \\
\hline
\texttt{FermionSPTLayersVerbose(SS, d)} & Compute and display layers up to dimension $d$ \\
\hline
\texttt{Spin12Factor(d, n)} & Get spin-1/2 factor for space group \\
\hline
\texttt{SptSetCohomology(CC, n)} & Compute $H^n$ from cochain complex \\
\hline
\end{tabular}
\end{center}

\subsection{Physics-to-Code Dictionary}

\begin{center}
\begin{tabular}{|l|l|}
\hline
\textbf{Physics Concept} & \textbf{GAP/SptSet Representation} \\
\hline
Symmetry group $G$ & GAP group object \\
\hline
Group cohomology $H^n(G, M)$ & \texttt{SptSetCohomology(CC, n)} \\
\hline
Coefficient module $M$ & \texttt{SptSetCoefficient*} objects \\
\hline
Spectral sequence $E_r^{p,q}$ & \texttt{SptSetSpecSeq*} objects \\
\hline
Fermion parity $(-1)^F$ & Homomorphism \texttt{auMap} to $\text{GL}(1,\mathbb{Z})$ \\
\hline
Spin structure & 2-cochain $w(g_1, g_2)$ \\
\hline
SPT classification & Output of \texttt{FermionSPTLayersVerbose} \\
\hline
\end{tabular}
\end{center}

\subsection{Where to Go Next}

\begin{enumerate}
    \item \textbf{Run examples}: Start with \texttt{examples/fspt\_2d\_ez.g}, then \texttt{fspt\_2d\_s12.g}
    \item \textbf{Read the original paper}: arXiv:2005.06572 explains the mathematical framework
    \item \textbf{Explore the code}: Begin with \texttt{lib/coefficient.gi} for simple coefficient systems
    \item \textbf{Check test files}: \texttt{tst/*.tst} files show expected inputs and outputs
\end{enumerate}

\section{The Spectral Sequence Behind \texttt{FermionSPTLayersVerbose}}
\label{sec:layers-pedagogical}

\subsection{Decoding the output}

The printed output of \texttt{FermionSPTLayersVerbose(SS, dim)} is best read as a compact window into the internal cache of a spectral sequence object \texttt{SS}.
Concretely:
\begin{itemize}
    \item The code fixes a \textbf{total degree} $p+q=\texttt{dim}+1$ and iterates over $p=1,2,\dots,\texttt{dim}+1$ with $q=\texttt{dim}+1-p$.
    \item A heading (``Bosonic'', ``Complex fermion'', ``Majorana'', ``p+ip'') is determined purely by the \textbf{$q$-value}:
    \begin{equation}
    q=0:\ \text{Bosonic},\quad q=1:\ \text{Complex fermion},\quad q=2:\ \text{Majorana},\quad q=3:\ \text{p+ip}.
    \end{equation}
    \item Each subsequent printed line under a fixed heading is \textbf{one page} $E_r^{p,q}$ at fixed $(p,q)$, for increasing $r$.
    Internally it calls \texttt{SptSetSpecSeqComponent(SS, r, p, q)} and then normalizes the result by \texttt{SptSetFpZModuleCanonicalForm} before printing.
    \item The number of printed lines is controlled by
    \begin{equation}
    r_{\max}=\max(q+2,\,p+1),\qquad r=2,3,\dots,r_{\max}.
    \end{equation}
    So ``why 3 lines?'' is a purely \textbf{algorithmic} question: e.g. for \texttt{dim=2}, ``Majorana'' means $(p,q)=(1,2)$ and hence $r_{\max}=4$, so it prints $E_2^{1,2},E_3^{1,2},E_4^{1,2}$.
\end{itemize}

Practical rule of thumb when reading the verbose output:
\begin{itemize}
    \item If the displayed group stops changing as $r$ increases, then this bidegree has stabilized (at least up to the printed cutoff), i.e. the relevant differentials are zero in this example.
    \item The printed group is shown in invariant-factor form: \texttt{<ZL-Module with torsions [ 2, 2, 2 ]>} means $\mathbb{Z}_2\times\mathbb{Z}_2\times\mathbb{Z}_2$.
\end{itemize}

Code pointers for this output:
\begin{itemize}
    \item \texttt{FermionSPTLayersVerbose}: \texttt{lib/ss\_fermion\_ez.gi}
    \item \texttt{SptSetSpecSeqComponent}: \texttt{lib/spectral\_sequence.gi}
    \item \texttt{SptSetFpZModuleCanonicalForm}: used to print abelian invariants
\end{itemize}

\subsection{Mathematical background: filtered complexes and the spectral sequence used here}

At the conceptual level, the spectral sequence in SptSet is used as a bookkeeping and computation device: it turns a difficult classification/cohomology problem into a sequence of successive approximations (``pages'') together with differentials (obstructions/identifications) that are applied step by step.

\paragraph{(1) Filtered cochain complexes and associated graded.}
Let $(C^{\bullet},d)$ be a cochain complex, and suppose we have a decreasing filtration by subcomplexes
\begin{equation}
\cdots \subseteq F^{p+1}C^{n} \subseteq F^{p}C^{n} \subseteq \cdots \subseteq C^{n}
\qquad\text{with}\qquad d(F^{p}C^{n})\subseteq F^{p}C^{n+1}.
\end{equation}
This induces a filtration on cohomology $H^{n}(C)$, and one defines the associated-graded pieces
\begin{equation}
\mathrm{gr}^{p}H^{n}(C) := F^{p}H^{n}(C)\big/ F^{p+1}H^{n}(C).
\end{equation}
A spectral sequence is a systematic way to compute these graded pieces (and track the extension problems needed to reconstruct $H^{n}(C)$).

\paragraph{(2) Bigrading and total degree.}
In most computational settings (including SptSet), one works with a bigraded object $C^{p,q}$ with a ``total'' grading $n=p+q$.
One canonical source is a double complex with two differentials of bidegrees $(1,0)$ and $(0,1)$ (up to sign conventions).
Filtering the total complex by $p$ produces a spectral sequence whose pages are indexed by $(r,p,q)$.

This explains the ubiquitous constraint ``$p+q=\texttt{dim}+1$'' in the user-facing output: the code is displaying a slice of fixed total degree, because that is the degree relevant for $d$-dimensional spatial phases (spacetime degree $d+1$).

\paragraph{(3) Pages and differentials.}
For a cohomological spectral sequence, the $r$-th differential has bidegree
\begin{equation}
d_{r}: E_{r}^{p,q}\longrightarrow E_{r}^{p+r,\,q-r+1}. \label{eq:ss_diff}
\end{equation}
The next page is obtained by taking cohomology with respect to this differential:
\begin{equation}
E_{r+1}^{p,q} \cong \ker(d_{r}:E_{r}^{p,q}\to E_{r}^{p+r,q-r+1})\big/\operatorname{im}(d_{r}:E_{r}^{p-r,q+r-1}\to E_{r}^{p,q}).
\end{equation}
In the package, the maps that ultimately determine these $d_r$ are installed when constructing \texttt{SS}; the function \texttt{SptSetSpecSeqComponent(SS, r, p, q)} then returns the resulting group $E_r^{p,q}$ (with caching).

\paragraph{(4) Convergence and the meaning of ``stabilization''.}
Under standard boundedness hypotheses (finite range of $p,q$ for fixed total degree), for each fixed $(p,q)$ the groups eventually stabilize:
\begin{equation}
E_{r}^{p,q}\cong E_{r+1}^{p,q}\cong\cdots\cong E_{\infty}^{p,q}\qquad\text{for }r\gg 0.
\end{equation}
The limiting groups give the associated graded pieces of the filtered target:
\begin{equation}
E_{\infty}^{p,q}\cong \mathrm{gr}^{p}H^{p+q}(C).
\end{equation}
This is the precise mathematical content behind the informal ``layer'' language.
The final classification group is then reconstructed from these graded pieces via extension data; there is no general guarantee that it is a direct sum.

\paragraph{(5) Why the project uses this viewpoint.}
Fermionic SPT/SET classifications typically come with multiple interacting pieces of data (often with different coefficient systems and obstruction relations).
Organizing them into bidegrees $(p,q)$ lets the package:
\begin{itemize}
    \item compute candidate pieces on an early page (often $E_2$),
    \item apply successive obstructions/identifications via differentials $d_r$, and
    \item expose the resulting stabilized pieces in a way that can be printed and inspected.
\end{itemize}
The headings ``Bosonic/Complex fermion/Majorana/p+ip'' are best understood as the package's naming convention for the low-$q$ sectors appearing in this filtration-based organization.

\subsection{Stacking and the group-extension problem (math $\to$ physics $\to$ code)}
\label{subsec:stacking-extension}

You now understand how a differential $d_r$ acts. Stacking is a different (but equally important) issue:
even after the spectral sequence stabilizes and you know the graded pieces $E_{\infty}^{p,q}\cong \mathrm{gr}^p H^{p+q}$, the final classification group is typically \emph{not} a direct sum of those pieces.
This is the classical \textbf{extension problem}.

\paragraph{(1) The mathematical extension problem.}
Let $H$ be the ``true'' classification group in total degree $n$ (in our application $n=d+1$). Suppose $H$ carries a filtration
\begin{equation}
H = F^{0}H \supseteq F^{1}H \supseteq \cdots \supseteq F^{n+1}H = 0,
\end{equation}
whose associated graded pieces are the stabilized spectral-sequence data
\begin{equation}
\mathrm{gr}^{p}H := F^{p}H/F^{p+1}H \cong E_{\infty}^{p,\,n-p}.
\end{equation}
Knowing all $\mathrm{gr}^{p}H$ does \emph{not} determine $H$ as a direct sum: for each $p$ there can be a nontrivial extension
\begin{equation}
0\to F^{p+1}H \to F^{p}H \to \mathrm{gr}^{p}H \to 0.
\end{equation}
Choosing representatives (a section) of $\mathrm{gr}^{p}H$ inside $F^{p}H$ is not canonical, and the failure of such a choice to respect addition is measured by a 2-cocycle valued in the lower filtration.

The simplest model is a short exact sequence of abelian groups
\begin{equation}
0 \to A \to E \to B \to 0.
\end{equation}
If we choose a set-theoretic section $s:B\to E$, then addition in $E$ is encoded by a \emph{twisting} term
\begin{equation}
s(b_1)+s(b_2) = s(b_1+b_2) + t(b_1,b_2),\qquad t(b_1,b_2)\in A.
\end{equation}
The map $t$ is exactly the extension 2-cocycle (it is zero iff the extension splits for this choice).

\paragraph{(2) What this means in the physics paper.}
In fermionic SPT classifications (e.g. arXiv:2310.19058), one often describes phases by \emph{layer data} (various cochains/cocycles) modulo equivalences, with obstruction relations between layers.
The spectral sequence packages these layers as the graded pieces, while the extension problem is the statement:
\begin{quote}
\emph{stacking two phases is not always ``add each layer separately''; there can be an extra correction term that lives in a lower layer.}
\end{quote}
In other words, the physical stacking law may contain cross-terms between the different pieces of data.
Those cross-terms are precisely what SptSet calls \textbf{add-twisters}.

\paragraph{(3) The translation to SptSet: what are $(p,q)$ for \texttt{SptSetInstallAddTwister}?}
Inside SptSet, a ``layered cochain representative'' of total degree $n$ is stored as a list
\begin{equation}
	exttt{layers} = [\,\texttt{layer}[0],\texttt{layer}[1],\dots,\texttt{layer}[n]\,],
\end{equation}
where \texttt{layer[p]} is an \emph{inhomogeneous $p$-cochain} (a GAP function of $p$ group elements) valued in the coefficient system indexed by
\begin{equation}
q = n-p.
\end{equation}
So the pair $(p,q)$ is exactly the bidegree label that you would also use for $E_r^{p,q}$.

Concretely, in the implementation of stacking (see \texttt{SptSetStack} in \texttt{lib/ss\_cochain.gi}) the code fixes the total degree
\begin{equation}
	exttt{deg} = n,
\end{equation}
loops over $p=0,1,\dots,n$, sets $q=n-p$, and then updates the $p$-layer by
\begin{equation}
	exttt{layers[p+1] := (c1.layer[p]) + (c2.layer[p]) + SS!.addTwister[p+1][q+1](c1.layers,c2.layers)}.
\end{equation}
This is exactly the extension-cocycle pattern $s(b_1)+s(b_2)=s(b_1+b_2)+t(b_1,b_2)$, except now there is one twister $t_{p,q}$ for each bidegree.

Therefore:
\begin{itemize}
    \item \texttt{SptSetInstallAddTwister(SS, p, q, A)} installs the twisting rule \emph{for the layer in bidegree $(p,q)$}.
    \item The function \texttt{A} is called as \texttt{A(layers1,layers2)} and must return a $p$-cochain (a GAP function of $p$ group elements) with values in the same coefficient module as the $(p,q)$ layer.
    \item It is natural (and common in the code) that \texttt{A} depends on \emph{other layers}, typically lower-$p$/higher-$q$ ones. This is the precise sense in which ``larger $q$ can affect smaller $q$'' when stacking: because $q=n-p$, a twister correcting a given $p$-layer may use data from smaller $p$ layers, which correspond to larger $q$.
\end{itemize}

\paragraph{(4) What is the relation of add-twisters to the page index $r$?}
The key point is that add-twisters are \emph{not} page-$r$ data.
The page index $r$ belongs to the spectral sequence computation of graded pieces (obstructions/identifications via $d_r$). In contrast, add-twisters describe how to realize the \emph{final classification group law} on chosen layered representatives.
So the installation API uses only $(p,q)$ (which layer is being corrected), but no $r$.

In practice, SptSet uses the same $(p,q)$ indexing for both:
\begin{itemize}
    \item differentials: \texttt{SptSetInstallCoboundary(SS, r, p, q, f)} (page-dependent),
    \item stacking twisters: \texttt{SptSetInstallAddTwister(SS, p, q, A)} (page-independent group law correction).
\end{itemize}

\paragraph{(5) How to read a concrete twister line.}
For example, the $(p,q)=(2,0)$ twister in the EZ constructor has the schematic form
\begin{equation}
	exttt{A(l1,l2) := \{g1,g2\} -> 1/2 * n11(g1) * n12(g2)},
\end{equation}
where the code first extracts \texttt{n11 := l1[1+1]} and \texttt{n12 := l2[1+1]}.
This should be read as:
\begin{itemize}
    \item we are stacking two total-degree-$n=2$ representatives (so layers correspond to $(p,q)=(0,2),(1,1),(2,0)$),
    \item the correction to the $(2,0)$ layer depends bilinearly on the two $(1,1)$ layers,
    \item the factor $\frac12$ indicates that the $(2,0)$ layer is $U(1)$-valued (additive notation), while the $(1,1)$ layers are $\mathbb{Z}_2$-valued.
\end{itemize}
This is precisely the ``cross-term'' in the stacking law that encodes a nontrivial extension between layers.

\section{Physical Inputs and Current Implementation Status}
\label{sec:physical-inputs-status}

This section answers a practical question that comes up immediately when trying to generalize the package:
\begin{quote}
\emph{What are the true ``physical inputs'' to the computation, what is currently implemented in code (differentials and stacking/group-extension rules), what is missing, and what additional code is needed for full $3{+}1$D and $4{+}1$D fermionic SPT (FSPT) classifications?}
\end{quote}

\subsection{What the physics inputs really are (at the cochain level)}

In this codebase, the physics does \emph{not} enter as a Hamiltonian or field theory; it enters as \textbf{cochain-level rules} that define:
\begin{enumerate}
    \item a space(-time) symmetry group $G$ (in practice a GAP group in a form suitable for resolutions),
    \item the coefficient systems (``spectrum'') that appear in the layered description,
    \item the obstruction/differential maps between layers (encoded as cochain formulas that induce $d_r$ on pages), and
    \item the stacking law (group law) on cocycle data, including the \textbf{twisting terms} that implement nontrivial group extensions between layers.
\end{enumerate}

Concretely, for the current fermion SPT pipelines the user provides:
\begin{itemize}
    \item \textbf{The symmetry group $G$}: e.g. a crystallographic space group turned into a Pcp group \texttt{SG1}.
    \item \textbf{A resolution $R$ of $G$}: computed by HAP (e.g. \texttt{ResolutionAlmostCrystalGroup}).
    \item \textbf{A parity/orientation homomorphism} $\rho: G\to \mathrm{GL}(1,\mathbb{Z})\cong\{\pm1\}$ (named \texttt{auMap} in code), used to define twisted coefficients and the function
    \begin{equation}
    s(g)=\frac{1-(g^{\rho})_{11}}{2}\in\{0,1\}.
    \end{equation}
    In the code the mapping named \texttt{auMap} (here denoted $\rho$) records whether a group element acts anti-unitarily (e.g. time reversal). Concretely, if $g^{\rho}=[[ -1 ]]$ then $s(g)=1$ and $g$ is anti-unitary; if $g^{\rho}=[[1]]$ then $s(g)=0$ and $g$ is unitary. This matches the physics paper's notation $s_1(g)$ in arXiv:2310.19058 (the indicator of anti-unitary elements).
    \item \textbf{Optionally, a spin-structure 2-cochain} $w\in C^2(G,\mathbb{Z}_2)$ (named \texttt{w} in code) for the spin-$\tfrac12$ variant \texttt{FermionSPTSpecSeq}.
    \begin{quotation}
    \noindent\textbf{Remark (physics meaning of \texttt{w}):} The input $w(g_1,g_2)\in\mathbb{Z}_2$ is used to encode the \emph{projective} (spin-$\tfrac12$) multiplication law of symmetry operators acting on fermions. In typical conventions one has
    \begin{equation}
    U(g_1)\,U(g_2)=(-1)^{w(g_1,g_2)}\,U(g_1g_2),
    \end{equation}
    so $w$ is a representative 2-cocycle describing a central extension/lift of the spatial symmetry into a double cover (Spin/Pin, depending on whether orientation-reversing elements are present).

    	\textbf{Relation to the paper's $\omega_2$:} in many physics conventions the group-extension cocycle $\omega_2\in Z^2(G,\mathbb{Z}_2^f)$ plays exactly this role (up to choice of representative and possible coboundary redefinitions), so the package's \texttt{w} should be viewed as the concrete cochain representative of that extension data.

    	\textbf{EZ vs spin-$\tfrac12$:} the ``EZ'' constructor \texttt{FermionEZSPTSpecSeq} does not take \texttt{w} as an input and corresponds to the case where this spin/projective extension data is taken to be trivial in the computation. For crystallographic space groups with spin-$\tfrac12$ degrees of freedom, a nontrivial $w$ is typically determined by the choice of spin structure/double cover (in this codebase often via helpers like \texttt{Spin12Factor}).
    \end{quotation}
\end{itemize}

Then the remaining ``physics'' is hard-coded in the package as cochain formulas for (i) coboundary/differential corrections and (ii) stacking twisters.

\subsection{What is implemented}

\paragraph{A. Differential/obstruction input currently present in code.}

At the level of implementation, the ``differentials'' are installed by calls of the form
\begin{equation}
	\texttt{SptSetInstallCoboundary(SS, r, p, q, f)}
\end{equation}
inside the constructors \texttt{FermionEZSPTSpecSeq} and \texttt{FermionSPTSpecSeq}.
These functions $f$ are cochain-level recipes that the spectral-sequence engine uses to build induced maps $d_r$, which implements Eq.\eqref{eq:ss_diff} exactly.

However, the generic engine in \texttt{lib/spectral\_sequence.gi} only implements the induced differentials for:
\begin{itemize}
    \item $r=2$ (the primary nontrivial page-to-page differential), and
    \item $r=3$ (built in a special way by solving for representatives using the $r=2$ data).
\end{itemize}
For $r\ge 4$ the engine prints ``not implemented'' and returns \texttt{fail}.

This is the most important limitation for $3{+}1$D and $4{+}1$D work: for total degree $p+q=4$ or $5$, one generically needs higher-page differentials (at least potentially $d_4$ and beyond).

\paragraph{B. Stacking law / group-extension input currently present in code.}

The stacking rule is implemented at the cochain level by the operation \texttt{SptSetStack} in \texttt{lib/ss\_cochain.gi}:
it adds each layer and then applies an additional correction term
\begin{equation}
	\texttt{SS!.addTwister[p+1][q+1](layers1, layers2)}
\end{equation}
for $p>0$.
These \textbf{add-twister} functions are installed in the fermion constructors using
\begin{equation}
	\texttt{SptSetInstallAddTwister(SS, p, q, A)}.
\end{equation}
Physically, these are exactly the ``stacking rules'' that encode nontrivial group extensions between layers (the group law on cocycle representatives is not a direct product).

In the current code:
\begin{itemize}
    \item The $2{+}1$D twisters include several explicit formulas and one table-driven correction term.
    \item For $(3{+}1)$D the code explicitly marks ``place holders for twisters'' and installs many of them as \texttt{ZeroCocycle@}, i.e. currently absent.
\end{itemize}

\paragraph{C. How equivalence / extension problems are handled.}

The package does not only print pages $E_r^{p,q}$.
It also has a notion of an actual \emph{classification group} as cochain-data modulo equivalences, implemented by:
\begin{itemize}
    \item \texttt{SptSetSpecSeqCochain} and stacking (\texttt{lib/ss\_cochain.gi}),
    \item \texttt{SptSetSpecSeqClassRep} (\texttt{lib/ss\_class.gi}), which stores a representative cochain and defines the induced group law.
\end{itemize}

This layer-aware class machinery is exactly where ``group extension'' issues live computationally.
But there is a critical current limitation: in \texttt{lib/ss\_class.gi} the routine \texttt{PartialPurify@} aborts with
\begin{quote}
	\texttt{Error("Purification at r>2 is not implemented.");}
\end{quote}
So the equivalence/purification logic is only implemented in a way that relies on $r=2$ data.
For higher-dimensional problems where $r\ge 3$ and higher-page structure is essential, the current ``class'' machinery must be extended.

\paragraph{D. The status of \texttt{lib/fermion\_ez.gd} and \texttt{lib/fermion\_ez.gi}.}

These two files are currently \textbf{empty} in this workspace, even though they are loaded by \texttt{init.g} and \texttt{read.g}.
At the moment they therefore contribute no declarations and no implementation.
In practice, the actual ``fermion EZ'' functionality lives in \texttt{lib/ss\_fermion\_ez.gd}/\texttt{.gi}.

Interpretation:
\begin{itemize}
    \item Most likely \texttt{fermion\_ez.gd/gi} are placeholders from an earlier refactor (intended to hold a higher-level interface or shared formulas),
    \item or a planned split where the ``physics formulas'' would live outside the spectral-sequence constructor.
\end{itemize}

\paragraph{E. What \texttt{lib/data/AddTwister2D.dat} is supposed to do (and its current state here).}

The table \texttt{AddTwister2DTable@} is loaded in \texttt{lib/ext\_data.gi} by
\begin{equation}
	\texttt{LoadExtDataFile@(9, "AddTwister2D.dat")}.
\end{equation}
It is then queried via \texttt{ExtData@(AddTwister2DTable@, ...)} inside the $(p,q)=(3,0)$ add-twister term (see \texttt{lib/ss\_fermion\_ez.gi} and \texttt{lib/ss\_fermion.gi}).

Physically, this is part of the \textbf{stacking/group-law correction} at total degree $3$ (the bosonic/top layer in $2$D runs), i.e. it encodes a discrete correction term that is inconvenient to express purely as a simple polynomial in cup products.

In \emph{this} workspace, \texttt{lib/data/AddTwister2D.dat} exists but is \textbf{empty}.
As a consequence:
\begin{itemize}
    \item \texttt{LoadExtDataFile@} loads an empty table,
    \item \texttt{ExtData@} will almost always fall back to returning $0$,
    \item therefore the table-driven twister contribution is effectively \textbf{disabled}.
\end{itemize}
This means that, strictly speaking, the current outputs may be missing part of the intended stacking law in $2$D unless the data file is populated.

\subsection{Is the code "ready" for full $3{+}1$D and $4{+}1$D if you add physics cochain formulas?}

Not yet.
Adding the correct cochain-level formulas for higher-dimensional differentials and extension rules is necessary, but the current code also lacks the \textbf{infrastructure} to use them in full generality.

More precisely:
\begin{itemize}
    \item \textbf{Higher-page differentials are not supported.} The spectral-sequence engine only implements induced maps for $r=2$ and a special $r=3$ path. For $3{+}1$D and $4{+}1$D one should expect that $d_4$ (and potentially higher) can matter.
    \item \textbf{Purification/equivalence logic is incomplete beyond $r=2$.} The class machinery currently errors for $r>2$. This blocks a robust computation of the final classification group (including extension issues) once higher-page structure is required.
    \item \textbf{Many $(3{+}1)$D add-twister terms are explicit placeholders.} Even if you know the physical stacking corrections, they must be installed as actual \texttt{addTwister} functions; currently several are set to zero.
    \item \textbf{Data-driven twisters need a generation pipeline.} The $2$D table file is empty here, and an analogous file \texttt{AddTwister2DEZ.dat} is referenced by \texttt{lib/data/.gitignore} but not present. For higher dimensions one must decide whether twisters are formulaic (cup products) or require precomputed lookup tables, and provide the corresponding files and loaders.
\end{itemize}

\paragraph{What you would need to add (code, not just formulas).}
To make the package able to attack full $3{+}1$D and $4{+}1$D classifications in a way consistent with the current architecture, the missing pieces are roughly:
\begin{enumerate}
    \item \textbf{General $d_r$ support for $r\ge 4$} in \texttt{lib/spectral\_sequence.gi} (extend \texttt{SptSetSpecSeqBuildDerivative} / \texttt{SptSetSpecSeqBuildDerivative2} so it can build induced $d_r$ from installed cochain-level recipes, rather than hard-coding only $r=2,3$).
    \item \textbf{General purification for $r>2$} in \texttt{lib/ss\_class.gi} (extend \texttt{PartialPurify@} so it can trivialize elements using $d_r$ data for higher $r$, not only $r=2$).
    \item \textbf{Complete the $(3{+}1)$D and $(4{+}1)$D add-twister table/functions} (implement the nontrivial stacking corrections for the relevant $(p,q)$ bidegrees at total degrees $4$ and $5$).
    \item \textbf{Possibly extend the spectrum (available $q$-sectors).} The current fermion pipelines build coefficients for $q\le 3$.
    For $4{+}1$D, depending on the intended physics model (e.g. whether additional invertible fermionic building blocks appear as independent layers), one may need to add further coefficient objects and corresponding differential/twister data.
    \item \textbf{Validation tooling.} Once the above is implemented, it becomes crucial to add targeted tests (in \texttt{tst/}) and/or reproducible example scripts to verify known classifications in low dimensions before trusting higher-dimensional outputs.
\end{enumerate}

In short: the current codebase already has a promising architecture (filtered cochains, stacking with twisters, and a spectral-sequence engine), but it is not yet a ``plug in cochain formulas and get all dimensions'' system.
For $3{+}1$D and $4{+}1$D, you should expect to write additional core-support code (especially higher $d_r$ and purification) in addition to supplying the physics formulas.

\section{How to Read \texttt{lib/ss\_fermion\_ez.gi}: \texttt{FermionEZSPTSpecSeq} and Cochain-Level Physics Input}
\label{sec:ss-fermion-ez-walkthrough}

This section is a detailed reading guide for the file \texttt{lib/ss\_fermion\_ez.gi}.
The goal is not to comment on every line, but to explain the \emph{repeated patterns} so that you can translate any line in that file into precise math/physics statements.

\subsection{What the file is doing at a high level (one mental picture)}

The file contains two conceptually different blocks:
\begin{enumerate}
    \item \textbf{A spectral-sequence constructor} \texttt{FermionEZSPTSpecSeq(R, auMap)} that builds a spectral sequence object \texttt{SS} and installs the physics inputs (cochain-level differential formulas and stacking rules).
    \item \textbf{Printing helpers} such as \texttt{FermionEZSPTLayersVerbose} and \texttt{FermionSPTLayersVerbose} that query the spectral sequence object and display selected components $E_r^{p,q}$.
\end{enumerate}

So when you read the file, keep this workflow in mind:
\begin{equation}
\boxed{\texttt{R},\ \texttt{auMap}}\ \xrightarrow{\ \texttt{FermionEZSPTSpecSeq}\ }\ \boxed{\texttt{SS}}\ \xrightarrow{\ \texttt{FermionSPTLayersVerbose}\ }\ \boxed{\text{printed }E_r^{p,q}\text{ groups}}.
\end{equation}

\subsection{GAP syntax: why it starts with \texttt{InstallMethod} and what is being defined}

The file begins with
\begin{equation}
	\texttt{InstallMethod(FermionEZSPTSpecSeq, ... , function(R, auMap) ... end)}.
\end{equation}
In GAP, an \emph{operation} (here \texttt{FermionEZSPTSpecSeq}) can have multiple implementations (\emph{methods}) selected by the runtime types of its arguments.
The list
\begin{equation}
	\texttt{[IsHapResolution, IsGeneralMapping]}
\end{equation}
declares that this method is used when the first input is a HAP resolution and the second input is a mapping (the parity/orientation map \texttt{auMap}).

Inside the method body, you will frequently see these GAP idioms:
\begin{itemize}
    \item \texttt{local x, y, z;}: declares local variables.
    \item Lists are 1-indexed: \texttt{spectrum[1]} is the first entry.
    \item Anonymous functions:
    \begin{itemize}
        \item \texttt{g -> expr} is a one-argument function.
        \item \texttt{\{g1, g2, g3\} -> expr} is a multi-argument shorthand.
        \item \texttt{function(a,b) ... end} is the long form.
    \end{itemize}
    \item A returned function can close over variables defined outside it (a \emph{closure}); this is used heavily to build cochain formulas depending on \texttt{spectrum}, \texttt{s}, etc.
\end{itemize}

\subsection{Where the physics goes: coefficients, differentials, stacking}

Within \texttt{function(R, auMap) ... end}, there are three key phases.

\paragraph{Phase 1: choose coefficients (the ``spectrum'').}
The code constructs a list \texttt{spectrum} whose entries are coefficient systems used for each $q$-sector.
In the EZ constructor you see
\begin{equation}
	\texttt{U(1),}\ \mathbb{Z}_2,\ \mathbb{Z}_2,\ \mathbb{Z}\quad\text{(up to indexing conventions).}
\end{equation}
Mathematically, this is specifying which coefficient module appears in each bidegree $E_r^{p,q}$.

\paragraph{Phase 2: install cochain-level differential formulas.}
Calls of the form
\begin{equation}
	\texttt{SptSetInstallCoboundary(SS, r, p, q, f)}
\end{equation}
store a cochain-level recipe $f$ into \texttt{SS}.
Later, the spectral-sequence engine uses those recipes to build the induced maps $d_r$ on pages.

\paragraph{Phase 3: install stacking (group law) and extension data.}
Calls of the form
\begin{equation}
	\texttt{SptSetInstallAddTwister(SS, p, q, A)}
\end{equation}
install the twisting term in the stacking law.
This is the implementation of the ``group extension problem'' at the cocycle-representative level: stacking is \emph{not} always componentwise addition; it may require extra correction cochains.

\subsection{\texttt{SptSetInstallCoboundary(SS, r, p, q, f)}}
\label{subsec:callback-interface}

This subsection explains the single most important programming pattern in SptSet: \emph{higher-order} GAP functions.
Once you understand it, essentially all installed differentials and twisters become straightforward.

\paragraph{(1) \texttt{n2} and \texttt{dn2} are functions (cochains).}
GAP treats functions as \emph{first-class values}: you can store a function in a variable, pass it as an argument, return it from another function, and call it later.

In SptSet, an inhomogeneous $p$-cochain is represented \emph{literally} by a GAP function of $p$ group elements. For example, a $2$-cochain is a function
\begin{equation}
n_2: G\times G\to A,\qquad (g_1,g_2)\mapsto n_2(g_1,g_2),
\end{equation}
and in code it is something you can call like \texttt{n2(g1,g2)}.

So when you see
\begin{verbatim}
function(n2, dn2)
    return function(g1,g2,g3,g4)
        ... n2(g1,g2) ... dn2(g1,g2,g3) ...
    end;
end
\end{verbatim}
the variables \texttt{n2} and \texttt{dn2} are ordinary function inputs, except that their \emph{values happen to be cochain-functions}.

\paragraph{(2) Why is there a ``function returning a function''?}
The installed object is a \emph{builder}:
\begin{equation}
	\texttt{f}\;:\; (n_p,\, dn_p)\longmapsto \bigl(d_r(n_p)\bigr),
\end{equation}
where $d_r(n_p)$ is itself a cochain of the appropriate degree.
Programming-wise:
\begin{itemize}
    \item the outer function is called \emph{once per input cochain} (or per generator representative),
    \item the returned inner function is the actual cochain $G^{k}\to A$ that SptSet will evaluate on tuples $(g_1,\dots,g_k)$.
\end{itemize}
This is exactly the same pattern as in ordinary mathematics: given a cochain $n$, you define a new cochain $\Phi(n)$ by writing an explicit formula in terms of $n$ and group elements.

\paragraph{(3) How does SptSet call the callback you installed?}
The method \texttt{SptSetInstallCoboundary(ss,r,p,q,f)} stores your builder \texttt{f} into the table
\begin{equation}
	\texttt{ss!.bdry[r+1][p+1][q+1] := f.}
\end{equation}
There are two important call sites:
\begin{itemize}
    \item \textbf{When computing the coboundary of a \emph{spectral-sequence cochain}} (see \texttt{SptSetSpecSeqCoboundarySL}):
    SptSet first forms the usual inhomogeneous coboundary
    \begin{equation}
    	\texttt{da := InhomoCoboundary@(coeff, a)}
    \end{equation}
    and then evaluates your builder as
    \begin{equation}
    	\texttt{SS!.bdry[r+1][p+1][q+1](a, da)}.
    \end{equation}
    In other words, in this pathway the second argument really is the coboundary $da$.

    \item \textbf{When building the induced map $d_r$ on pages} (see \texttt{SptSetSpecSeqBuildDerivative}):
    for $r=2$ the engine currently calls
    \begin{equation}
    	\texttt{SS!.bdry[2+1][p+1][q+1](np\_, ZeroCocycle@)}
    \end{equation}
    where \texttt{np\_} is a representative inhomogeneous cocycle obtained from a module generator.
    This is why many formulas are written with guards like \texttt{if dn2 <> ZeroCocycle@ then ... fi;}.
\end{itemize}

\paragraph{(4) What is \texttt{dn2} mathematically?}
When it is provided as a genuine function (i.e. not \texttt{ZeroCocycle@}), \texttt{dn2} is the inhomogeneous coboundary of \texttt{n2} with the appropriate coefficient action:
\begin{equation}
dn_2 = \delta n_2 \in C^3(G,A).
\end{equation}
In the common $\mathbb{Z}_2$-valued, trivial-action situations, the explicit inhomogeneous formula is
\begin{equation}
(\delta n_2)(g_1,g_2,g_3)
= ~ ^{g_1} n_2(g_2,g_3) - n_2(g_1 g_2, g_3) + n_2(g_1, g_2 g_3) - n_2(g_1,g_2)\quad (\mathrm{mod}\ 2).
\end{equation}
(The general version with a nontrivial $G$-action is implemented by \texttt{InhomoCoboundary@} via \texttt{coeff!.gAction}.)

\subsubsection{A fully explicit example: the \texttt{d2: E2\textasciicircum\{2,1\} $\to$ E2\textasciicircum\{4,0\}} block}
\label{subsec:d2-21-example}

Consider the code fragment
\begin{verbatim}
val := 1/2 * (n2(g1, g2) * n2(g3, g4) mod 2);
if dn2 <> ZeroCocycle@ then
    val := val + 1/2 * ((n2(g1*g2*g3, g4) * dn2(g1, g2, g3)
        + n2(g1, g2*g3*g4) * dn2(g2, g3, g4)) mod 2);
    val := val + 1/2 * (dn2(g1, g2, g3*g4) * dn2(g1*g2, g3, g4) mod 2);
    val := val - 1/4 * (dn2(g1, g2, g3) * (1 - dn2(g1, g2, g3*g4)) mod 2);
fi;
\end{verbatim}

Here \texttt{n2} is a $\mathbb{Z}_2$-valued 2-cochain, and (when supplied) \texttt{dn2} is its 3-coboundary.
The returned cochain is a $U(1)$-valued 4-cochain written in additive notation with fractions $\frac12,\frac14$.
It comes from the following obstruction function 
\begin{align}
\mathrm{d}_{s_1} \nu _{3} &=\mathcal{O}_{4}[ n_{2}]( 01234)\nonumber\\\nonumber
&=( -1)^{[ \omega _{2} \cup n_{2} +n_{2} \cup n_{2} +n_{2} \cup _{1}\mathrm{d} n_{2} +\mathrm{d}( s_{1} \cup n_{2} +n_{2} \cup _{2}\mathrm{d} n_{2})]( 01234)}\\\nonumber
&\quad\times ( -1)^{\omega _{2}( 013)\mathrm{d} n_{2}( 1234) +\mathrm{d} n_{2}( 0124)\mathrm{d} n_{2}( 0234)}\\
&\quad\times ( -i)^{\mathrm{d} n_{2}( 0123)[ 1-\mathrm{d} n_{2}( 0124)] (\mathrm{mod}\ 2)}
.
\end{align}
Notice that the coboundary $\frac{1}{2} \mathrm{d}( s_{1} \cup n_{2} +n_{2} \cup _{2}\mathrm{d} n_{2})$ does not enter the formula, which is not important for classification, but useful for stacking.
Please keep this in mind for later debug.

\subsection{A worked translation example (small)}

Consider the installed coboundary formula
\begin{equation}
	\texttt{SptSetInstallCoboundary(SS, 2, 1, 2, function(n1,dn1) return \{g1,g2,g3\}-> s(g1)*n1(g2)*n1(g3); end)}.
\end{equation}
Read this as:
\begin{itemize}
    \item Input layer: a $\mathbb{Z}_2$-valued 1-cochain $n_1\in C^{1}(G,\mathbb{Z}_2)$.
    \item Output: a degree-3 cochain $\Phi\in C^{3}(G,\cdot)$ defined by
    \begin{equation}
    \Phi(g_1,g_2,g_3)= s(g_1)\,n_1(g_2)\,n_1(g_3),
    \end{equation}
    where $s(g)\in\{0,1\}$ is computed from \texttt{auMap} via $s(g)=\frac{1-(g^{\texttt{auMap}})_{11}}{2}$.
\end{itemize}
This is exactly the style used throughout the file: the physics is encoded as explicit polynomial-like expressions in inhomogeneous cochains and cup-$i$ products.

\subsection{A worked translation example (stacking / extension data)}

Now look at an add-twister installation:
\begin{equation}
	\texttt{SptSetInstallAddTwister(SS, 2, 0, function(l1,l2) ... return \{g1,g2\} -> 1/2*n11(g1)*n12(g2); end)}.
\end{equation}
Here the inputs \texttt{l1} and \texttt{l2} are not group elements; they are the \emph{layer lists} of two total-degree cochain representatives being stacked.

\paragraph{(1) What exactly are \texttt{l1} and \texttt{l2}?}
In the implementation, a ``phase/cochain representative'' is stored as an object of type \texttt{IsSptSetSpecSeqCochainRep} with a field \texttt{.layers}.
For a fixed total degree \texttt{deg} (denoted $n$ in the math discussion), the field
\begin{equation}
	\texttt{c!.layers}
\end{equation}
is a GAP list of length \texttt{deg+1} indexed by $p=0,1,\dots,\texttt{deg}$.
The entry
\begin{equation}
	\texttt{c!.layers[p+1]}
\end{equation}
is an inhomogeneous $p$-cochain (i.e. a GAP function of $p$ group elements), valued in the coefficient sector labeled by
\begin{equation}
q = \texttt{deg} - p.
\end{equation}

So in \texttt{SptSetInstallAddTwister(SS,p,q,A)}, the function \texttt{A} is called with
\begin{equation}
	\texttt{l1 = c1!.layers},\qquad \texttt{l2 = c2!.layers}
\end{equation}
coming from the stacking operator \texttt{SptSetStack}.
Concretely, \texttt{SptSetStack} computes (for each $p$ with $q=\texttt{deg}-p$)
\begin{verbatim}
layers[p+1] := AddInhomoCochain@(c1!.layers[p+1], c2!.layers[p+1]);
layers[p+1] := AddInhomoCochain@(layers[p+1],
    SS!.addTwister[p+1][q+1](c1!.layers, c2!.layers));
\end{verbatim}
This shows that \texttt{l1} and \texttt{l2} are \emph{not} “lists of twists”; they are the raw layer data of the two inputs being stacked.
The twister \texttt{A(l1,l2)} is the cross-term that corrects naive layerwise addition.

\paragraph{(2) Why does \texttt{l1} look like ``Majorana then complex fermion''?}
Because for a fixed total degree \texttt{deg}, each index $p$ automatically determines $q=\texttt{deg}-p$.
In the fermion EZ spectrum, the code uses the naming convention
\begin{equation}
q=0:\ \text{bosonic }U(1),\quad q=1:\ \text{complex fermion }\mathbb{Z}_2,\quad q=2:\ \text{Majorana }\mathbb{Z}_2,\quad q=3:\ p{+}ip.
\end{equation}
So, for example, when \texttt{deg=3} we have
\begin{equation}
(p,q)=(1,2)\ \text{is a Majorana 1-cochain},\qquad (p,q)=(2,1)\ \text{is a complex-fermion 2-cochain}.
\end{equation}
That is why the code snippet you pasted for the $(p,q)=(3,0)$ twister reads
\begin{equation}
	\texttt{n11 := l1[1+1];\ \ n21 := l1[2+1];}
\end{equation}
and comments them as “Majorana 1-cochain” and “complex fermion 2-cochain”: it is simply reading the $p=1$ and $p=2$ layers of a total-degree-$3$ object.

\paragraph{(3) What is stored in a layer entry?}
Each entry \texttt{l1[p+1]} is either
\begin{itemize}
    \item \texttt{ZeroCocycle@}, the distinguished sentinel meaning the zero cochain (a constant-zero function), or
    \item a genuine cochain function, e.g. \texttt{n11(g)} for $p=1$ or \texttt{n21(g1,g2)} for $p=2$.
\end{itemize}
This is why add-twister code frequently checks whether two layers are \texttt{ZeroCocycle@} before constructing more complicated expressions.

\paragraph{(4) Back to the simple $(2,0)$ example.}
The code extracts
\begin{equation}
	\texttt{n11 := l1[2];\quad n12 := l2[2]}
\end{equation}
which you should read as ``take the $p=1$ layer of the first cochain and the $p=1$ layer of the second cochain'' (index shift because GAP lists start at 1).
Then the returned function
\begin{equation}
(g_1,g_2)\mapsto \tfrac12\,n_{1}^{(1)}(g_1)\,n_{1}^{(2)}(g_2)
\end{equation}
is exactly a correction term added during stacking.
This is the computational incarnation of a nontrivial extension in the final classification group.

\paragraph{(5) How to read the $(3,0)$ bosonic twister you pasted.}
In the code block
\begin{equation}
	\texttt{SptSetInstallAddTwister(ss, 3, 0, function(l1,l2) ... end)}
\end{equation}
the twister corrects the \emph{bosonic 3-cochain layer} $(p,q)=(3,0)$ when stacking two total-degree-$3$ representatives.
It constructs intermediate cochains from lower layers:
\begin{itemize}
    \item \texttt{n11 := l1[1+1]}, \texttt{n12 := l2[1+1]} are the two Majorana 1-cochains (bidegree $(1,2)$) from the two inputs.
    \item \texttt{n21 := l1[2+1]}, \texttt{n22 := l2[2+1]} are the two complex-fermion 2-cochains (bidegree $(2,1)$) from the two inputs.
    \item \texttt{dn21}, \texttt{dn22} are not taken from \texttt{l1/l2}; they are \emph{constructed} from \texttt{n11/n12} using the known relation (a specific $d_2$-type formula used in the physics input).
    \item \texttt{m2} and \texttt{N2} are auxiliary 2-cochains built from the lower layers; these are the cross-terms that make the stacking law non-product.
\end{itemize}
Finally it returns a 3-cochain
\begin{equation}
	\texttt{\{g1,g2,g3\} -> (1/2*c3(g1,g2,g3) + t3(g1,g2,g3))},
\end{equation}
which is exactly the correction term $t(b_1,b_2)$ (in the extension-problem sense) added into the $(3,0)$ layer.
The \texttt{ExtData@(...)} piece \texttt{t3} is a table-driven correction term (encoding a case distinction that is awkward to express as a simple polynomial in cup products).

\subsection{How to decode the rest of the file yourself (a checklist)}

For any line in \texttt{lib/ss\_fermion\_ez.gi}, you can translate it by following this mechanical procedure:
\begin{enumerate}
    \item Identify whether you are looking at a \textbf{coboundary installation} (\texttt{SptSetInstallCoboundary}) or a \textbf{stacking twister} (\texttt{SptSetInstallAddTwister}).
    \item Read off the indices $(r,p,q)$ (for coboundaries) or $(p,q)$ (for twisters); those indices tell you which bidegree the formula belongs to.
    \item Interpret each symbol:
    \begin{itemize}
        \item \texttt{n1}, \texttt{n2}, \texttt{n3} are inhomogeneous cochains in degrees $1,2,3$ respectively.
        \item \texttt{dn1}, \texttt{dn2}, \texttt{dn3} represent their coboundaries.
        \item \texttt{s(g)} encodes the parity/orientation data from \texttt{auMap}.
        \item \texttt{Cup0@, Cup1@, Cup2@} are cup-$i$ products producing higher cochains.
        \item \texttt{AddInhomoCochain@} is addition in the target coefficient group.
    \end{itemize}
    \item Translate the returned function \texttt{\{g1,...\} -> expr} into an explicit formula $\Phi(g_1,\dots)=\text{expr}$.
\end{enumerate}

Once you can apply this recipe, every line in \texttt{FermionEZSPTSpecSeq} becomes readable: it is a list of explicit cochain formulas organized by bidegree.

\section{How to Read \texttt{lib/ss\_vanilla.gi}: The General AHSS Engine}
\label{sec:ss-vanilla-structure}

This section explains the generic Atiyah--Hirzebruch spectral sequence (AHSS) engine implemented in \texttt{lib/ss\_vanilla.gi}.
This file is the down-layer that \texttt{lib/ss\_fermion\_ez.gi} and \texttt{lib/ss\_fermion.gi} build upon.
The role of \texttt{ss\_vanilla.gi} is to turn a resolution plus a list of coefficient systems into a working spectral sequence object whose pages and differentials can be queried.
All physics input lives in the installed coboundary and stacking rules.
The infrastructure for pages, kernels, images, and induced maps is here.

\subsection{High-level mathematical picture}

Fix a symmetry group $G$ and a resolution $R$ of $\mathbb{Z}$ as a $\mathbb{Z}G$-module.
Choose a spectrum of coefficient systems $\{A_q\}$ indexed by $q\ge 0$.
The AHSS in this code is organized as
\begin{equation}
E_1^{p,q} = C^p(G, A_q),
\end{equation}
with differentials
\begin{equation}
d_r: E_r^{p,q} \to E_r^{p+r, q-r+1}.
\end{equation}
The implemented recursion for $r\ge 2$ is
\begin{equation}
E_r^{p,q} = \ker(d_{r-1}: E_{r-1}^{p,q} \to E_{r-1}^{p+r-1, q-r+2})\;\big/\;\operatorname{im}(d_{r-1}: E_{r-1}^{p-r+1, q+r-2} \to E_{r-1}^{p,q}).
\end{equation}
The ``physics'' formulas in the fermion files determine the $d_r$ by specifying explicit cochain-level maps.
The algebraic bookkeeping above is universal and lives in \texttt{ss\_vanilla.gi}.

\subsection{The spectral-sequence object and its fields}

The representation \texttt{IsSptSetSpecSeqVanillaRep} extends \texttt{IsSptSetSpecSeqRep} with the fields
\begin{equation}
\texttt{resolution},\ \texttt{brMap},\ \texttt{spectrum},\ \texttt{bdry},\ \texttt{addTwister},\ \texttt{cochainTypes},\ \texttt{classTypes}.
\end{equation}
These are initialized in \texttt{SptSetSpecSeqVanilla(resolution, spectrum)}.
Each entry has a precise mathematical role.

\paragraph{\texttt{resolution}.}
This is the HAP resolution $R$ of the group $G$.
It provides the chain modules $R_p$ and differentials needed to construct cochains and cocycles.

\paragraph{\texttt{brMap}.}
This is the ``bar resolution map'' produced by \texttt{SptSetBarResolutionMap(resolution)}.
It converts between HAP resolution cocycles and inhomogeneous cochains on $G$.
The physics formulas are written in the inhomogeneous language.
This map is the bridge that allows $d_r$ to be computed from explicit inhomogeneous expressions.

\paragraph{\texttt{spectrum}.}
This is a list $\{A_q\}$ of coefficient systems.
For the fermion EZ case this list is built in \texttt{FermionEZSPTSpecSeq} and looks like $U(1)$, $\mathbb{Z}_2$, $\mathbb{Z}_2$, $\mathbb{Z}$ up to indexing conventions.
The $q$-index is always accessed as \texttt{spectrum[q+1]} due to GAP's 1-based lists.

\paragraph{\texttt{bdry}.}
This is a 3D array of installed cochain-level coboundary builders.
It is populated by \texttt{SptSetInstallCoboundary(ss, r, p, q, f)}.
The function \texttt{f} encodes the physics formula for the $d_r$-type obstruction at bidegree $(p,q)$.

\paragraph{\texttt{addTwister}.}
This table stores the stacking corrections installed by \texttt{SptSetInstallAddTwister(ss, p, q, A)}.
It is not used to build $E_r$ pages directly.
It is crucial for the group law on SPT phases.
It is later accessed by \texttt{SptSetStack} when adding two cocycle representatives.

\paragraph{\texttt{cochainTypes} and \texttt{classTypes}.}
These caches store GAP types for cochain representatives and their equivalence classes at a fixed total degree.
They are created by \texttt{SptSetSpecSeqCochainType} and \texttt{SptSetSpecSeqClassType} so that objects carry the correct spectral-sequence metadata.

\subsection{Building pages: \texttt{SptSetSpecSeqBuildComponent}}

The function \texttt{SptSetSpecSeqBuildComponent(ss, r, p, q)} implements the formula for $E_r^{p,q}$.
There are three key branches.

\paragraph{(1) Bounds and missing spectrum.}
If $p<0$ or $q<0$, or if \texttt{spectrum[q+1]} is not bound, the component is defined to be the zero module.
This is the code's way of truncating the AHSS outside the supported bidegrees.

\paragraph{(2) First page.}
For $r=1$ the code returns
\begin{equation}
E_1^{p,q} = C^p(G, A_q),
\end{equation}
implemented by \texttt{SptSetCochainModule(resolution, p, spectrum[q+1])}.
This is the inhomogeneous cochain module, not yet modded by any differential.

\paragraph{(3) Higher pages.}
For $r\ge 2$ the code builds the homology of the previous page.
It first constructs the incoming differential
\begin{equation}
\phi = d_{r-1}: E_{r-1}^{p-r+1,\,q+r-2} \to E_{r-1}^{p,q}
\end{equation}
and the outgoing differential
\begin{equation}
\psi = d_{r-1}: E_{r-1}^{p,q} \to E_{r-1}^{p+r-1,\,q-r+2}.
\end{equation}
Then it returns \texttt{SptSetHomologyModule(phi, psi)}, which is the algebraic module $\ker \psi / \operatorname{im} \phi$.
This is the abstract definition of the $E_r$ page encoded in code.

\subsection{The tilde-page: \texttt{SptSetSpecSeqBuildComponent2}}

The variant \texttt{SptSetSpecSeqBuildComponent2} builds
\begin{equation}
\widetilde{E}_r^{p,q} = \operatorname{coker}\bigl(d_{r-1}: E_{r-1}^{p-r+1,\,q+r-2} \to E_{r-1}^{p,q}\bigr).
\end{equation}
It uses \texttt{SptSetCokernelModule(phi)} rather than a full homology module.
This ``tilde page'' is used internally as the target for \texttt{SptSetSpecSeqBuildDerivative2}.
This avoids double quotienting when the target is already known to be a quotient.
In practice, this affects how the image representatives are projected when assembling the matrix for $d_r$.
The underlying math is that when you already know the target bidegree, you can represent it as a cokernel rather than as a kernel modulo an image.
The cokernel of a linear map $f:A\to B$ is
\begin{equation}
\operatorname{coker}(f)=B/\operatorname{im}(f).
\end{equation}
This is the natural dual concept to the kernel, and it measures how far $f$ is from being surjective.
In the AHSS, the map at page $r-1$ is
\begin{equation}
d_{r-1}: E_{r-1}^{p,q} \to E_{r-1}^{p+r-1,q-r+2}.
\end{equation}
The cokernel of this specific map is precisely the object defined in the code as
\begin{equation}
\widetilde{E}_r^{p+r,q-r+1} = \operatorname{coker}(d_{r-1}).
\end{equation}
Accordingly, there is an exact sequence
\begin{equation}
E_{r-1}^{p,q} \xrightarrow{d_{r-1}} E_{r-1}^{p+r-1,q-r+2} \longrightarrow \widetilde{E}_r^{p+r,q-r+1} \to 0,
\end{equation}
where the second arrow is the quotient map sending an element of $E_{r-1}^{p+r-1,q-r+2}$ to its class modulo the image of $d_{r-1}$.
This is the reason the target here is the tilde page, not the full $E_r$ page.
The true page is
\begin{equation}
E_r^{p+r,q-r+1} = \ker\bigl(d_{r-1}: E_{r-1}^{p+r-1,q-r+2} \to E_{r-1}^{p+2r-2,q-2r+3}\bigr)\big/\operatorname{im}(d_{r-1}),
\end{equation}
so it is a subquotient of the same middle module.
In other words, $\widetilde{E}_r$ forgets the kernel condition on the outgoing differential and keeps only the quotient by the incoming image.
That is exactly why it is convenient as a computational target for building $d_r$.
Computationally, representing the target as a cokernel means you only quotient by \emph{one} image, not by a kernel and an image.
This is why \texttt{SptSetSpecSeqBuildDerivative2} maps into \texttt{SptSetSpecSeqBuildComponent2} rather than the full homology module.
For users, the takeaway is that the ``tilde page'' is not a new mathematical object; it is a practical representation that stabilizes the computation of $d_r$ by avoiding redundant quotienting.
When you interpret outputs, remember that $
\widetilde{E}_r^{p,q}$ is only an internal model for the same $E_r^{p,q}$ target.

\subsection{Building differentials: \texttt{SptSetSpecSeqBuildDerivative}}

The core algorithm for $d_r$ lives in \texttt{SptSetSpecSeqBuildDerivative} and follows a strict chain of conversions.
The input is the module $M = E_r^{p,q}$ and the output is a linear map to $N = E_r^{p+r,q-r+1}$.
The algorithm computes the images of the generators of $M$ one by one.

\paragraph{Step 1: choose a cocycle representative.}
Each generator of $M$ is mapped to an inhomogeneous cocycle using
\begin{equation}
\texttt{np\_ := SptSetMapToBarCocycle(brMap, p, spectrum[q+1], generator)}.
\end{equation}
The result \texttt{np\_} is a $p$-cochain on $G$ representing that generator.

\paragraph{Step 2: apply the cochain-level spectral-sequence coboundary.}
The function
\begin{equation}
\texttt{dnp := SptSetSpecSeqCoboundarySL(ss, p+q, p, np\_)}
\end{equation}
returns a total-degree $(p+q+1)$ cochain whose layers encode the installed physics formula for $d_r$.
Here the total degree $p+q$ is the physical dimension of the phase being classified.
The code then enforces the cocycle condition on the next layer by setting
\begin{equation}
\texttt{dnp!.layers[p+2] := ZeroCocycle@}.
\end{equation}

\paragraph{Step 3: convert to a class and purify.}
The cochain \texttt{dnp} is turned into a class object and purified on page $r$ by
\begin{equation}
\texttt{cl\_dnp := SptSetSpecSeqClassFromCochainNC(dnp);}
\end{equation}
\begin{equation}
\texttt{PartialPurifySSClass@(cl\_dnp, r, p+r-1);}
\end{equation}
This step kills lower-page ambiguities so that the remaining leading layer is exactly the obstruction living at bidegree $(p+r, q-r+1)$.
The assertion
\begin{equation}
\texttt{LeadingLayer(cl\_dnp) = p+r-1}
\end{equation}
is the algebraic statement that the obstruction vanishes on previous pages.
Mathematically, purification is the operation of replacing a cocycle by an equivalent representative after quotienting by the images of lower-page differentials.
Equivalence here means ``differ by a $d_{r-1}$-exact term'' in the spectral-sequence sense.
In physics language, this is a gauge choice for the layered cochain representative: you can shift lower layers by coboundaries without changing the phase.
The routine \texttt{PartialPurifySSClass@} explicitly subtracts such lower-page contributions so that the first nonzero layer appears at $p+r-1$.
This matters because $d_r$ is defined on the quotient $\ker d_{r-1}/\operatorname{im} d_{r-1}$.
Without purification, the ``leading layer'' you extract could still contain pieces that are trivial on the $E_r$ page.
Purification guarantees that the extracted layer is a well-defined element of $E_r^{p+r,q-r+1}$ rather than a representative that depends on the previous-page gauge.
When you use or debug the code, a good check is to verify that after purification the earlier layers vanish and the leading layer index matches $p+r-1$.

\paragraph{Step 4: extract the obstruction layer.}
The actual $d_r$ image is taken as the negative of the leading layer,
\begin{equation}
\texttt{opr\_ := -cl\_dnp!.cochain!.layers[p+r+1]}
\end{equation}
followed by conversion back to a bar cocycle and a linear map.
This yields one column of the matrix for $d_r$.
Repeating for all generators gives \texttt{SptSetZLMapByImages(M, N, fA)}.

\paragraph{Why this matches physics.}
The cochain-level formulas installed in \texttt{ss\_fermion\_ez.gi} and \texttt{ss\_fermion.gi} are exactly the obstruction functions appearing in the AHSS for fermionic SPTs.
The vanilla engine is agnostic about their origin.
It only enforces the algebraic constraints of a spectral sequence.
This separation is deliberate.
The physics tells the engine what the cochain formula is.
The engine handles kernels, quotients, and representative choices.

\subsection{The role of \texttt{SptSetInstallCoboundary} and \texttt{SptSetInstallAddTwister}}

The interfaces
\begin{equation}
\texttt{SptSetInstallCoboundary(ss, r, p, q, f)}
\end{equation}
and
\begin{equation}
\texttt{SptSetInstallAddTwister(ss, p, q, A)}
\end{equation}
are defined here because they are universal infrastructure.
They are not fermion-specific logic.
The coboundary table is queried by \texttt{SptSetSpecSeqCoboundarySL} when building differentials.
The add-twister table is queried by \texttt{SptSetStack} when two class representatives are added.
Together they encapsulate the full ``physics input'' needed by the AHSS machinery.

\subsection{Practical decoding checklist for \texttt{ss\_vanilla.gi}}

When reading the file line-by-line, keep the following map in mind.

\begin{itemize}
    \item \textbf{\texttt{SptSetSpecSeqVanilla}} creates a spectral sequence object with empty tables and cached page slots.
    \item \textbf{\texttt{SptSetSpecSeqBuildComponent}} encodes $E_r^{p,q} = \ker d_{r-1} / \operatorname{im} d_{r-1}$.
    \item \textbf{\texttt{SptSetSpecSeqBuildComponent2}} encodes the cokernel version needed for internal targets.
    \item \textbf{\texttt{SptSetSpecSeqBuildDerivative}} converts generators into inhomogeneous cocycles, applies the physics coboundary, purifies, and extracts the obstruction layer.
    \item \textbf{\texttt{SptSetSpecSeqBuildDerivative2}} is the same algorithm, but targets the tilde-page to avoid extra quotienting.
    \item \textbf{\texttt{SptSetSpecSeqCochainType} / \texttt{SptSetSpecSeqClassType}} only manage GAP typing and caching, but are essential for consistent object behavior.
\end{itemize}

If you can track the flow of a single generator through Steps 1--4 above, then you have the core of the AHSS engine in your hands.
At that point, the ``down layer'' is fully demystified.
The physics sections in \texttt{ss\_fermion\_ez.gi} and \texttt{ss\_fermion.gi} become direct inputs to a standard algebraic machine.

\section{Group Extension and Stacking Infrastructure: \texttt{ss\_module}, \texttt{ss\_class}, \texttt{ss\_cochain}}
\label{sec:ss-extension-stacking}

This section explains the group-extension (stacking) layer in the \texttt{group\_ext} branch.
The three files form a strict dependency chain.
The bottom layer is \texttt{ss\_cochain.gi}, which defines layered cochains and the stacking law.
The middle layer is \texttt{ss\_class.gi}, which defines equivalence classes and purification.
The top layer is \texttt{ss\_module.gi}, which assembles $E_\infty$ pieces into the final classification via extensions.
If you want to debug stacking or extension bugs, \texttt{ss\_module.gi} is the first file to read, but it only makes sense once you understand the two lower layers.

\subsection{Where the extension problem comes from (math and physics)}

For a fixed total degree $n=p+q$, the AHSS produces graded pieces $E_\infty^{p,q}$.
These pieces are not yet the final classification group.
They are the associated graded of a filtration
\begin{equation}
0 = F_{-1} \subset F_0 \subset F_1 \subset \cdots \subset F_n = \mathcal{C}_n,
\end{equation}
with
\begin{equation}
F_p/F_{p-1} \cong E_\infty^{p,n-p}.
\end{equation}
The ``group extension problem'' is the problem of reconstructing $\mathcal{C}_n$ from the graded pieces.
In physics, this is the statement that stacking SPT phases is not always a direct product of layers.
Instead, lower-layer data can twist the addition law of higher layers.
In code, this twisting is captured by \texttt{addTwister} (at the cochain level) and by the extension logic in \texttt{ss\_module.gi} (at the module level).

\subsection{Bottom layer: \texttt{ss\_cochain.gi} (stacking cochain representatives)}

The file \texttt{ss\_cochain.gi} defines the data model for a spectral-sequence cochain and the stacking rule.
An object of type \texttt{IsSptSetSpecSeqCochainRep} is a list of layers
\begin{equation}
\texttt{layers[p+1]} \in C^p(G, A_{n-p}),
\end{equation}
for total degree $n$.
The constructor \texttt{SptSetSpecSeqCochain(ss, deg, layers)} just stores these layers with the correct type.
The helper \texttt{SptSetSpecSeqCochainZero(ss, deg)} initializes all layers as \texttt{ZeroCocycle@}.

The key operation is stacking.
The method
\begin{equation}
\texttt{SptSetStack(c1, c2)}
\end{equation}
performs layerwise addition and then applies a twist:
\begin{equation}
\texttt{layers[p+1] := layers[p+1] + addTwister[p+1][q+1](l1,l2)}.
\end{equation}
This implements the physical group law, not the naive direct sum of layers.
The twister is only applied for $p>0$, which matches the code path that forbids twisting the $p=0$ layer.
The inverse \texttt{SptSetInverse} is implemented by negating layers and compensating the twist term at each $p$.
This means a wrong \texttt{addTwister} will break the group law in a nonlocal way.
When debugging, first verify that stacking then inversion returns the zero cochain for low-dimensional test inputs.

The other key routine here is \texttt{SptSetSpecSeqCoboundarySL}.
It takes a single layer $a\in C^p(G,A_{n-p})$, computes its ordinary inhomogeneous coboundary $\delta a$, and then appends the spectral-sequence correction layers built from \texttt{bdry}.
This is how a single cochain input is lifted into a full layered coboundary representative.
The class and module layers below rely on this function, so bugs in \texttt{bdry} or its indexing will propagate to extension results.

\subsection{Middle layer: \texttt{ss\_class.gi} (classes and purification)}

The class layer upgrades cochains to equivalence classes under the stacking law and coboundaries.
An object of type \texttt{IsSptSetSpecSeqClassRep} stores a single field \texttt{cochain}.
Addition of classes is defined by stacking the underlying cochains.
The zero class is the zero cochain of the same total degree.
The additive inverse is defined using \texttt{SptSetInverse} from the cochain layer.

The most important function is \texttt{SptSetPurifySpecSeqClass}.
It iteratively removes lower-layer coboundaries so that the first nonzero layer appears at the correct position.
This is a computational version of choosing a representative in the quotient $\ker d_{r-1}/\operatorname{im} d_{r-1}$.
In practice it uses two subroutines.
The function \texttt{PartialPurify@} removes contributions that are trivial on a finite page $r$ by solving a linear equation for a preimage under $d_r$.
The function \texttt{PartialPurifyCoboundary@} removes trivial ordinary coboundaries ($r=1$) by solving a cocycle equation.
Both of these steps modify the cochain in place and also return the ``correction'' cochain that was stacked in.
This is why class purification is essential for extension computations: it gives a stable leading layer that can be mapped into modules.

If you see inconsistent extension results, check that \texttt{SptSetPurifySpecSeqClass} is actually converging.
The code currently throws an error if purification for $r>2$ is requested, which is a known limitation.
This limitation means higher-page extension logic is not reliable beyond $r=2$ without further implementation.

\subsection{Top layer: \texttt{ss\_module.gi} (assembling the extension)}

The file \texttt{ss\_module.gi} is the top of the stacking.
It constructs explicit abelian groups from the $E_\infty$ pages and glues adjacent layers via extension data.
This is where the group extension problem is solved in code.

\paragraph{Data model.}
An object of type \texttt{IsSptSetSpecSeqModuleRep} stores
\begin{equation}
\texttt{specSeq},\ \texttt{deg},\ \texttt{pRange},\ \texttt{basis\_classes},\ \texttt{components},\ \texttt{vector\_embedings}.
\end{equation}
Here \texttt{deg} is the total degree $n$ and \texttt{pRange} is the list of $p$ values being assembled.
Each \texttt{components[p]} is the stabilized module $E_\infty^{p,n-p}$.
The array \texttt{basis\_classes} stores class representatives corresponding to the module generators.
The matrices \texttt{vector\_embedings[p]} map component vectors into the combined ambient lattice.

\paragraph{From a single page to a module: \texttt{SptSetSpecSeqComponentEx}.}
This function takes one bidegree $(p,q)$ and returns a SpecSeq module with \texttt{pRange = [p]}.
It calls \texttt{SptSetSpecSeqComponentInf} to obtain $E_\infty^{p,q}$ and then canonicalizes its generators.
Each generator is lifted to a class using \texttt{SptSetSpecSeqClassFromLevelCocycle}.
These classes are stored as a basis for later extension glueing.
This step already depends on the cochain and class layers.

\paragraph{Vector--class conversion.}
The conversion \texttt{SptSetSpecSeqModuleVectorToClass(M,a)} takes an integer vector and returns a class by summing basis classes.
The reverse conversion is harder.
The function \texttt{SptSetSpecSeqModuleClassToLeadingVector} purifies the class, reads off its leading layer, and then maps that layer into the module basis.
The full conversion \texttt{SptSetSpecSeqModuleClassToVector} iteratively peels off leading-layer contributions across the \texttt{pRange}.
This loop is delicate.
If purification fails or the leading layer is misidentified, the resulting vector will be wrong and extension logic will break.

\paragraph{The extension constructor: \texttt{SptSetSpecSeqModuleExtension}.}
This function glues two adjacent ranges, say $pRange(M_1)=\{p_0,\dots,p_f-1\}$ and $pRange(M_2)=\{p_f\}$.
It builds a new module with generators and relations encoded by three matrices \texttt{Emat}, \texttt{Pmat}, and \texttt{Rmat}.
If either input is zero, the extension is trivial.
Otherwise the function creates block matrices that contain the generators and relations of both modules.

The key extension data enters through torsion relations in $M_1$.
For a torsion generator with relation coefficient $t_j$, the code computes
\begin{equation}
c_{j} = \texttt{SptSetSpecSeqModuleVectorToClass(M1, t_j, e_j)},
\end{equation}
then maps it into the next layer using
\begin{equation}
v_j = \texttt{SptSetSpecSeqModuleClassToLeadingVector(M2, c_j)}.
\end{equation}
This vector $v_j$ is inserted into the off-diagonal part of \texttt{Rmat}.
That off-diagonal block is the algebraic record of the extension class.
Physically, it encodes how stacking a torsion element in the lower layer produces a nontrivial element in the next layer.
If you suspect a stacking bug, inspect these $v_j$ values first.
They directly reflect the add-twister data and the purification logic.

\paragraph{Assembling the full classification: \texttt{SptSetSpecSeqResult}.}
The function \texttt{SptSetSpecSeqResult(ss, deg, pRange)} iterates through $pRange$ and repeatedly applies \texttt{SptSetSpecSeqModuleExtension}.
This constructs the filtered group $\mathcal{C}_n$ from its graded pieces.
If you change the stacking law or the purification logic, this is the function that will reveal the consequences.

\paragraph{Induced maps between extensions.}
The method \texttt{SptSetMapInducedByGroupHomomorphism} maps generators from one SpecSeq module to another via a group homomorphism.
It converts a module vector to a class, maps the underlying cochain, and then converts back to a vector.
This path depends on all three layers and is a good stress test for correctness.

\subsection{Practical debugging checklist for stacking and extension}

Check \texttt{addTwister} first.
An incorrect twister makes the group law non-associative and will corrupt all extensions.
Verify that \texttt{SptSetStack} is associative on a small test case.

Check purification next.
Use a small group and verify that \texttt{SptSetPurifySpecSeqClass} produces the same leading layer after repeated calls.
If the leading layer changes across calls, your coboundary data or index ranges are inconsistent.

Check extension for torsion.
For a torsion generator $t_j e_j$, confirm that the computed $v_j$ is zero in trivial cases and nonzero only when physics predicts an extension.
The function \texttt{SptSetSpecSeqModuleClassToLeadingVector} is the usual culprit if this fails.

Finally, validate the full classification.
Run \texttt{SptSetSpecSeqResult} for a toy input where the classification is known to be a direct product.
If the output is not a direct product, the extension logic or the twister input is wrong.


\section{Debugging: ``Top Layer Is Not a Cocycle'' Assertion Failure}
\label{sec:debug-top-layer-not-cocycle}

This section documents a specific assertion failure encountered when running 2D wallpaper group examples, along with troubleshooting steps and the underlying cause.

\subsection{The error}

When running \texttt{examples/fspt\_2d\_ez\_ext.g}, wallpaper group 12 fails with:
\begin{verbatim}
12
Layers:
[ <ZL-Module with torsions [ 2 ]>, <ZL-Module with torsions [ 2 ]>,
  <ZL-Module with torsions [ 2, 2, 4 ]> ]
Error, Assertion failure: ASSERTION FAILURE: top layer is not a cocycle in
  Assert( -1, ForAll( cp, IsInt ),
    "ASSERTION FAILURE: top layer is not a cocycle" );
 at /opt/gap-4.15.1/pkg/SptSet/lib/ss_class.gi:74 called from
SptSetPurifySpecSeqClass( cl ); at ss_module.gi:80
\end{verbatim}

\subsection{Call stack analysis}

The error propagates through:
\begin{enumerate}
    \item \texttt{SptSetSpecSeqResult(SS, 3, [1,2,3])} in \texttt{fspt\_2d\_ez\_ext.g:22}.
    \item \texttt{SptSetSpecSeqModuleExtension(M, Epq)} in \texttt{ss\_module.gi:236}.
    \item \texttt{SptSetSpecSeqModuleClassToLeadingVector(M2, cjn)} in \texttt{ss\_module.gi:188}.
    \item \texttt{SptSetPurifySpecSeqClass(cl)} in \texttt{ss\_module.gi:80}.
    \item The assertion at \texttt{ss\_class.gi:74}.
\end{enumerate}

\subsection{What the assertion checks}

At line 74 of \texttt{ss\_class.gi}, the code performs:
\begin{verbatim}
cp := SptSetMapFromBarCocycle(brMap, p, SS!.spectrum[q+1], cp_);
Assert(-1, ForAll(cp, IsInt),
       "ASSERTION FAILURE: top layer is not a cocycle");
\end{verbatim}
The function \texttt{SptSetMapFromBarCocycle} converts an inhomogeneous cochain \texttt{cp\_} into a coordinate vector \texttt{cp} with respect to the HAP resolution basis.

For $U(1)$ coefficients (used at $q=0$, the bosonic top layer), this involves a Bockstein map.
The Bockstein of a true cocycle is an integer vector.
If \texttt{cp\_} is \emph{not} a cocycle (i.e., $\delta(\texttt{cp\_}) \neq 0$), the Bockstein returns fractional values, and the assertion fails.

\subsection{Root cause: cochain is not closed}

The assertion failure means that during purification, a cochain layer \texttt{layers[p+1]} that should represent a cocycle in $H^p(G, A_{n-p})$ is actually not closed under the coboundary.
This can happen for several reasons:
\begin{enumerate}
    \item \textbf{Incorrect \texttt{addTwister} data.}
    If the twister at $(p,q)=(3,0)$ returns wrong values, stacking two valid cochains produces a non-cocycle.
    \item \textbf{Bug in \texttt{SptSetStack} or \texttt{SptSetInverse}.}
    The stacking and inversion logic at the cochain level may introduce numerical errors.
    \item \textbf{Missing or incorrect differential installation.}
    If \texttt{SptSetInstallCoboundary} for certain $(p,q,r)$ is wrong, the coboundary computation is corrupted.
    \item \textbf{Group-specific resolution issues.}
    Some wallpaper groups have more complex resolutions.
    Group 12 (p6, hexagonal) has a non-trivial point group $C_6$, which may expose edge cases.
\end{enumerate}

\subsection{Debugging procedure}

\paragraph{Step 1: Isolate the failing group.}
Create a debug script that runs only group 12:
\begin{verbatim}
LoadPackage("SptSet");
it := 12;
SG := SpaceGroupBBNWZ(2, it);
fSG := IsomorphismPcpGroup(SG);
SG1 := Image(fSG);
R := ResolutionAlmostCrystalGroup(SG1, 6);
gs := GeneratorsOfGroup(SG);

f := GroupHomomorphismByImagesNC(SG1, GL(1, Integers),
  List(gs, x -> Image(fSG, x)),
  List(gs, x -> [[DeterminantMat(x)]]));

SS := FermionEZSPTSpecSeq(R, f);
layers := FermionSPTLayers(SS, 2);
Display(layers);
\end{verbatim}
This should succeed and produce the layers shown in the error.
The failure occurs later, when \texttt{SptSetSpecSeqResult} calls the extension machinery.

\paragraph{Step 2: Check individual layer classes.}
Use the pattern from \texttt{debug/p31m\_ez\_ext.g}:
\begin{verbatim}
E30inf := SptSetSpecSeqComponentInf(SS, 3, 0);
SptSetFpZModuleCanonicalForm(E30inf);
Display(E30inf);

# Check layer 1 generators
n := SptSetNumberOfGenerators(layers[1]);
for i in [1..n] do
  v1 := layers[1]!.generators[i];
  cl1 := SptSetSpecSeqClassFromLevelCocycle(SS, 3, 1, v1);
  cl2 := cl1 + cl1;  # Stack with itself
  SptSetPurifySpecSeqClass(cl2);  # May fail here
  Display(LeadingLayer(cl2));
od;
\end{verbatim}
This will pinpoint which generator causes the failure.

\paragraph{Step 3: Inspect the cochain before failure.}
Before calling \texttt{SptSetPurifySpecSeqClass}, examine the cochain:
\begin{verbatim}
cl2 := cl1 + cl1;
Display(cl2!.cochain!.layers);
\end{verbatim}
Check if layer values are reasonable (should be small integers or simple fractions from $U(1)$).

\paragraph{Step 4: Trace the \texttt{addTwister} contribution.}
The $(p,q)=(3,0)$ twister is loaded from \texttt{AddTwister2D.dat}.
Verify that:
\begin{itemize}
    \item The file \texttt{lib/data/AddTwister2D.dat} is not empty.
    \item The file was copied from \texttt{FSPT Stacking Theta.txt} correctly.
    \item The 9-parameter lookup matches expectations.
\end{itemize}
In the GAP session:
\begin{verbatim}
Display(AddTwister2DTable@);  # Should be a non-empty list
ExtData@(AddTwister2DTable@, 0,0,0,0,0,0,0,0,0);  # Should be 0
ExtData@(AddTwister2DTable@, 1,0,1,0,0,1,1,0,0);  # Lookup a specific entry
\end{verbatim}

\paragraph{Step 5: Check coboundary installation.}
Verify that the coboundaries are installed correctly in \texttt{ss\_fermion\_ez.gi}:
\begin{verbatim}
Display(SS!.bdry);  # Inspect the installed coboundary functions
\end{verbatim}
Pay attention to $(p,q,r)$ combinations that involve the failing layer.

\paragraph{Step 6: Test stacking associativity.}
Create three simple cochains and verify $(a \oplus b) \oplus c = a \oplus (b \oplus c)$:
\begin{verbatim}
c1 := SptSetSpecSeqClassFromLevelCocycle(SS, 3, 1, v1);
c2 := SptSetSpecSeqClassFromLevelCocycle(SS, 3, 1, v1);
c3 := SptSetSpecSeqClassFromLevelCocycle(SS, 3, 1, v1);
left := (c1 + c2) + c3;
right := c1 + (c2 + c3);
# Compare left!.cochain!.layers with right!.cochain!.layers
\end{verbatim}
If these differ, the twister or stacking logic is broken.

\subsection{Known limitations and workarounds}

\paragraph{Purification only works for $r \leq 2$.}
The function \texttt{PartialPurify@} explicitly checks \texttt{Assert(0, r <= 2)} and will fail for higher-page differentials.
This means groups requiring $d_3$ or higher purification will fail.
For 2D examples, $r=2$ should suffice.

\paragraph{Some wallpaper groups may have bugs.}
Not all 17 wallpaper groups have been exhaustively tested.
Groups 11--17 involve more complex point groups (triangular, hexagonal lattices) and may expose edge cases in the resolution or twister data.

\paragraph{Workaround: skip failing groups.}
Modify the example loop to skip problematic groups:
\begin{verbatim}
for it in [2..17] do
  if it in [12] then continue; fi;  # Skip known failures
  # ... rest of loop
od;
\end{verbatim}
This allows testing other groups while investigating the failure.

\subsection{Summary}

The ``top layer is not a cocycle'' error indicates that a cochain that should be closed is not.
The most likely causes are:
\begin{enumerate}
    \item Incorrect or missing \texttt{AddTwister2D.dat} data.
    \item A bug in the stacking law for the specific group.
    \item An edge case in the resolution or coboundary for that wallpaper group.
\end{enumerate}
Follow the debugging steps above to isolate the failing component.
Start with the \texttt{addTwister} data, then check stacking associativity, and finally inspect the coboundary installation.


\appendix

\section{GAP Cheat Sheet}
\label{app:gap-cheat-sheet}

This appendix is a compact reminder of GAP syntax (and a few conventions used throughout SptSet).

\subsection{Core syntax}

\begin{verbatim}
# comments start with #
x := 5;            # assignment uses :=
x = 5;             # equality test uses =
x <> 5;            # inequality test uses <>

true; false;       # booleans
not b;             # boolean negation
b1 and b2;         # boolean and
b1 or b2;          # boolean or

# ';' prints values, ';;' suppresses printing
expr;
expr;;
\end{verbatim}

\subsection{Control flow}

\begin{verbatim}
if cond then
    ...
elif otherCond then
    ...
else
    ...
fi;

for x in list do
    ...
od;

while cond do
    ...
od;
\end{verbatim}

\subsection{Functions and calling conventions}

\begin{verbatim}
f := x -> x^2;                           # short lambda
f := function(x) return x^2; end;       # full syntax

# varargs (used heavily for inhomogeneous cochains)
g := function(arg...) return Length(arg); end;

# calling a function with a list of arguments
CallFuncList(f, [x]);
\end{verbatim}

\subsection{Lists, records, and indexing}

\begin{verbatim}
# lists
L := [1, 2, 3];
Length(L);
Add(L, 4);            # append
Remove(L, 2);         # remove and return L[2]

# list slicing (1-based indices!)
L{[2..3]};             # sublist [2,3]

# records (like dictionaries)
R := rec(a := 1, b := 2);
R.a;                   # component access

# SptSet frequently uses '!.' access on objects/records
# SS!.bdry, coeff!.gAction, etc.
\end{verbatim}

\subsection{Help and inspection}

\begin{verbatim}
?Name;                 # help for a name
??pattern;             # search in help system
Display(obj);          # pretty print
\end{verbatim}


\subsection{What does \texttt{@} mean in SptSet code?}
\label{app:gap-at-sptset}

You will see identifiers like \texttt{ZeroCocycle@}, \texttt{Cup1@}, \texttt{AddInhomoCochain@} inside the SptSet source files. In the GAP package system this trailing \texttt{@} indicates a \emph{package-namespace-qualified} global.

Concretely:
\begin{itemize}
    \item Inside the package implementation files, the name is written with a trailing \texttt{@}, e.g. \texttt{ZeroCocycle@}.
    \item From a GAP session (outside the package source files), the same global can be referred to as \texttt{ZeroCocycle@SptSet}. Many scripts in \texttt{examples/} and \texttt{debug/} use this style.
\end{itemize}

This convention avoids name clashes with other packages and makes it visually obvious which helpers are provided by SptSet.

\section{SptSet Cheat Sheet}
\label{app:sptset-cheat-sheet}

This appendix is a compact reminder of the most frequently used SptSet workflow and its key objects.

\subsection{Load and run}

\begin{verbatim}
LoadPackage("SptSet");
\end{verbatim}

\subsection{Typical 2D workflow (space group -> resolution -> spectral sequence -> layers)}

\begin{verbatim}
# 1) choose a space group (from cryst)
SG := SpaceGroupBBNWZ(2, it);          # it = 1..17 for wallpaper groups

# 2) convert to a Pcp group (needed by HAP resolution algorithms)
fSG := IsomorphismPcpGroup(SG);
SG1 := Image(fSG);

# 3) compute a (HAP) resolution up to a dimension bound
R := ResolutionAlmostCrystalGroup(SG1, 6);

# 4) build the fermion parity / orientation map f: G_b -> GL(1,Z)
gs := GeneratorsOfGroup(SG);
f := GroupHomomorphismByImagesNC(
    SG1, GL(1, Integers),
    List(gs, x -> Image(fSG, x)),
    List(gs, x -> [[DeterminantMat(x)]])
);

# 5) construct a spectral sequence object
SS := FermionEZSPTSpecSeq(R, f);       # easy parity case

# 6) read off the classification layers in low dimension
FermionSPTLayersVerbose(SS, 2);
\end{verbatim}

\subsection{Core SptSet objects (how to recognize them)}

\begin{itemize}
    \item \texttt{R}: a HAP resolution (e.g. \texttt{ResolutionAlmostCrystalGroup}).
    \item \texttt{SS}: an SptSet spectral sequence object (a GAP object with internal components accessed as \texttt{SS!.\,something}).
    \item \texttt{cochains}: inhomogeneous cochains are represented as GAP functions of several group elements,
    e.g. \texttt{n2 := \{g1,g2\} -> ...} or \texttt{n2 := function(g1,g2) ... end}.
\end{itemize}

\subsection{Indexing conventions (why \texttt{+1} appears everywhere)}

GAP lists are 1-based, while the mathematics of spectral sequences is naturally indexed starting at 0.
Therefore the code frequently uses shifts like \texttt{p+1}, \texttt{q+1}, \texttt{r+1} when accessing arrays of pages/differentials, e.g.
\begin{verbatim}
SS!.bdry[r+1][p+1][q+1]( ... );
\end{verbatim}

\subsection{Inhomogeneous cochain helpers (most used)}

All of the following live in the SptSet namespace (see \S\ref{app:gap-at-sptset}).

\begin{verbatim}
ZeroCocycle@;                 # canonical "zero cochain" sentinel
InhomoCoboundary@(coeff, a);  # coboundary da of an inhomogeneous cochain a

AddInhomoCochain@(a, b);      # addition of cochains
NegativeInhomoCochain@(a);    # -a
ScaleInhomoCochain@(k, a);    # k*a

Cup0@(p, q, coeff, a, b);     # cup product a \smile b
Cup1@(p, q, coeff, a, b);     # cup-1 product
Cup2@(p, q, coeff, a, b);     # cup-2 product
\end{verbatim}


\subsection{Spectral sequence plumbing (when you develop/extend the package)}

The core installation points are:
\begin{verbatim}
SptSetInstallCoboundary(ss, r, p, q, formula);
SptSetInstallAddTwister(ss, p, q, twister);
\end{verbatim}
where \texttt{formula} and \texttt{twister} are usually given as anonymous functions returning inhomogeneous cochains.

\subsection{Worked 3+1D example: $G_f=\mathbb Z_2^f \times \mathbb Z_4^T$ from the S4 point-group input}
\label{sec:worked-s4-z2f-z4t}

This subsection records the explicit computational trace corresponding to \texttt{examples/fspt\_ez\_ext.g} at the point-group entry \texttt{S4}.
The matching debug script is \texttt{debug/debug\_s4\_ext.g}.
In this setup, the bosonic quotient is $G_b \simeq \mathbb Z_4$, the orientation character is nontrivial, and the 3+1D layer output is
\begin{equation}
\left(E_\infty^{1,3}, E_\infty^{2,2}, E_\infty^{3,1}, E_\infty^{4,0}\right)
= \left(\mathbb Z_2, 0, \mathbb Z_2, 0\right).
\end{equation}
Thus only the $p+ip$ layer and the complex-fermion layer survive at $E_\infty$.

\paragraph{Step 1: choose a $p+ip$ generator}
Let
\begin{equation}
x = (n_1, n_2, n_3, \nu_4) \in E_\infty^{1,3}
\end{equation}
be the basis class extracted by \texttt{SptSetSpecSeqComponentEx(ss, 1, 3)}.
The script prints all nonzero entries of the inhomogeneous cochains $n_1, n_2, n_3, \nu_4$ on indexed group elements of $G_b$.
This gives a concrete representative, not only the abstract module class.

\paragraph{Step 2: stack two copies}
Stacking is computed by
\begin{equation}
x+x
= \left(2n_1,\; 2n_2 + t_{22}(n_1,n_1),\; 2n_3 + t_{31}(n_2,n_2),\; \nu_4^{(2)}\right),
\end{equation}
with add-twisters $t_{pq}$ installed in \texttt{ss\_fermion\_ez.gi}.
In the rollback configuration used here, the $(2,2)$ twister is disabled, so
\begin{equation}
t_{22}(n_1,n_1)=0.
\end{equation}
The script still shows a nontrivial $(3,1)$ twister contribution in canonical coordinates:
\begin{equation}
t_{31}(x,x)\neq 0 \in E_\infty^{3,1}\cong \mathbb Z_2.
\end{equation}
This is the key source of the extension $2[p+ip]=[\text{complex fermion}]$.
At this point a common confusion is whether $2n_1$ should already be zero because the final layer is $\mathbb Z_2$.
The answer is no at the cochain level.
In the debug output, $N_1=2n_1$ is explicitly nonzero as a function on group elements.
However, its canonical coordinate in $E_\infty^{1,3}$ is already zero.
So $N_1$ is a nonzero coboundary representative of the zero class.

\paragraph{Step 3: purification}
The class $x+x$ initially has nonzero lower-layer representatives.
Therefore purification is necessary, and it is indeed performed on the $p+ip$ layer first.
The call \texttt{SptSetPurifySpecSeqClass(x2)} returns a correction cochain $b$ and updates $x_2\leftarrow x_2+\mathrm d b$.
In this S4 run, the script prints
\begin{equation}
b_0=-1,\qquad b_1\neq 0,\qquad b_2=b_3=0.
\end{equation}
So the first correction is a degree-0 piece that kills the nonzero-but-trivial $N_1$ layer.
After \texttt{SptSetPurifySpecSeqClass}, the $p=1,2$ layers are removed as coboundary-trivial data, while the $p=3$ component remains nontrivial:
\begin{equation}
x+x \xrightarrow{\text{purify}} y,\qquad y\in E_\infty^{3,1}\setminus\{0\}.
\end{equation}
The debug print verifies that the canonical coordinate of $N_3'(x+x)$ equals the canonical coordinate of the basis generator in $E_\infty^{3,1}$.
So the clean statement is
\begin{equation}
\text{cochain level: }2n_1\neq 0,
\qquad
\text{class level: }2[x]=0\in E_\infty^{1,3},
\end{equation}
and purification is exactly the step that implements this equivalence.

\paragraph{Step 4: extension module}
The module-extension trace prints
\begin{equation}
M_{12} \cong \mathbb Z_2,\qquad
M_{123} \cong \mathbb Z_4,\qquad
M_{1234} \cong \mathbb Z_4,
\end{equation}
with relation matrix collapsing to a single torsion-$4$ generator at the $(1,3)+(3,1)$ stage.
Equivalently,
\begin{equation}
2x = y,\qquad 2y = 0.
\end{equation}
Therefore the final interacting 3+1D classification block in this channel is
\begin{equation}
\mathbb Z_4,
\end{equation}
rather than $\mathbb Z_2 \times \mathbb Z_2$.

\paragraph{Interpretation}
The extension does not come from a surviving Majorana $E_\infty^{2,2}$ class (that layer is zero).
It comes from the nontrivial stacking rule that feeds the doubled $p+ip$ class into the complex-fermion layer, and then survives purification as the unique nontrivial class in $E_\infty^{3,1}$.


\section{Debug note: $p+ip$ extension mechanism in SptSet code}
\label{sec:debug_pip_extension}

This section documents the findings from debugging the $G_f = \mathbb{Z}_4^{f,T}$ ($G_b = \mathbb{Z}_2^T$, $\omega_2 = m_1^2$, $s_1 = m_1$) example in the SptSet code.
The notation follows \texttt{3+1D\_Stacking.tex}.

\subsection{The addTwister(2,2) puzzle}

In \texttt{3+1D\_Stacking.tex}, the $p+ip$ root state is $(n_1, n_2, n_3, \nu_4) = (m_1, 0, 0, 1)$.
When stacking two copies, the $\mathbb{Z}_2$-valued Majorana layer should be
\begin{equation}
N_2 = n_2 + n_2' + m_2 = 0 + 0 + s_1 \cup (n_1 \cup_1 n_1') = m_1^2 \mod 2,
\end{equation}
where $m_2 = s_1 \cup (n_1 \cup_1 n_1')$ is the addTwister(2,2) contribution.
This twister is essential: without it, $N_2 = 0$ and the $p+ip \to \mathrm{MC}$ extension would be trivial.

However, in the SptSet code, addTwister(2,2) is set to zero, yet the $p+ip \to \mathrm{MC}$ extension is \emph{still nontrivial} ($\mathbb{Z}_4$).
This seems contradictory.

\subsection{Resolution: $\mathbb{Z}$-valued cochain representation}

The key is that the SptSet code does \emph{not} store $\mathbb{Z}_2$-valued cochains.
Instead, all cochain layers are stored as $\mathbb{Z}$-valued functions (integer-valued).
The $\mathbb{Z}_2$ torsion is encoded in the \emph{bar resolution module relations}, not in the cochain values themselves.

Concretely, the $p+ip$ basis class in the code is a degree-4 cochain $(c_0, c_1, c_2, c_3, c_4)$ where each $c_p$ is an integer-valued function.
Its construction proceeds as follows:
\begin{enumerate}
\item Start with the level cocycle $c_1 = m_1$ (the $p+ip$ generator in $H^1(G_b, \mathbb{Z}_T)$).
\item Compute the spectral sequence coboundary $\mathrm{d}c_1$ using the installed differential formulas.
Since $m_1$ is a cocycle in $\mathbb{Z}_T$ coefficient, the ordinary coboundary $\mathrm{d}_{s_1} m_1 = 0$.
\item The higher-layer contributions from the differentials $d_2, d_3, \ldots$ at $(p,q) = (1,3)$ are computed.
For the $d_2$ formula $d_{2,1,3}(n_1) = \omega_2 \cdot n_1 + s_1 \cdot n_1 \cdot n_1 = m_1^3 + m_1^3 = 2m_1^3$, which equals $0 \mod 2$.
\item The purification procedure then ``solves'' the coboundary equations to fill in the tails $c_2, c_3, c_4$.
\end{enumerate}

The crucial result from the code: the $p=2$ layer of the $p+ip$ basis class is
\begin{equation}
c_2(T, T) = -1, \qquad c_2(\text{otherwise}) = 0.
\end{equation}
This is a $\mathbb{Z}$-valued function, \emph{not} reduced mod 2.
Physically, $c_2 \equiv -1 \equiv 1 \mod 2$, so it represents a nontrivial element in $H^2(G_b, \mathbb{Z}_2)$.
But as a $\mathbb{Z}$-valued function, $c_2$ carries ``half-integer information'' about the lift from $\mathbb{Z}_2$ to $\mathbb{Z}$.

\subsection{The stacking mechanism in $\mathbb{Z}$-lift representation}

When the code stacks two copies of the $p+ip$ class (before purification):
\begin{align}
N_1 &= c_1 + c_1 = 2m_1, \\
N_2 &= c_2 + c_2 + \underbrace{\mathrm{addTwister}(2,2)}_{= 0} = 2c_2 = -2.
\end{align}
Here $N_2(T,T) = -2$ as an integer.

Then the purification step removes $N_1 = 2m_1 = \mathrm{d}_{s_1} m_0$ (with $m_0 = -1$) by subtracting the spectral sequence coboundary of $m_0$.
This coboundary propagates to the $p=2$ layer through the differential formula:
\begin{equation}
d_{2,0,3}(m_0) = m_0 \cdot \omega_2 = (-1) \cdot m_1^2.
\end{equation}
Evaluating at $(T,T)$: $(-1) \cdot m_1^2(T,T) = (-1) \cdot 1 = -1$.

After purification:
\begin{equation}
N_2' = N_2 - d_{2,0,3}(m_0) = -2 - (-1) = -1.
\end{equation}
Wait --- more precisely, the coboundary of $m_0$ shifts $N_2$ by $-d_{2,0,3}(m_0) = +1 \cdot m_1^2 = +1$ at $(T,T)$, 
so $N_2' = -2 + (-1) = -3$.
(The exact sign depends on the code's convention for ``subtract coboundary.'')

The numerical result from the code confirms: $N_2'(T,T) = -3$.
In the bar resolution representation, this maps to the vector $[-3]$, which is odd and hence nontrivial in the $\mathbb{Z}_2$ module.

\paragraph{Comparison with the manual calculation.}
In the manual calculation, one works mod 2 throughout:
\begin{equation}
N_2 = 0 + 0 + m_2 = m_1^2 \quad (\text{mod } 2).
\end{equation}
In the code's $\mathbb{Z}$-lift calculation:
\begin{equation}
N_2' = -3 \equiv 1 \quad (\text{mod } 2).
\end{equation}
Both give the generator $m_1^2$ of $H^2(G_b, \mathbb{Z}_2) = \mathbb{Z}_2$.
The mechanisms are equivalent: the manual twister $m_2$ and the code's purification correction both produce the same $\mathbb{Z}_2$ cohomology class, just via different $\mathbb{Z}$-lifts.

This is why enabling addTwister(2,2) \emph{double-counts}: the purification already encodes the same information.

\subsection{The remaining bug: $p+ip$ breaks the MC $\to$ CF extension}

The SptSet code computes group extensions layer by layer.
Without $p+ip$, the chain MC $\to$ CF $\to$ Bosonic correctly gives $\mathbb{Z}_8$.
With $p+ip$, the chain $p+ip \to$ MC $\to$ CF $\to$ Bosonic gives $\mathbb{Z}_4 \times \mathbb{Z}_4$ instead of $\mathbb{Z}_{16}$.

The failure point: after extending $p+ip$ by MC to get $M_{12} = \mathbb{Z}_4$, the next step extends $M_{12}$ by CF.
The canonical generator of $M_{12}$ is a ``mixed'' class $g = -(\text{p+ip class}) + (\text{MC class})$.
The extension requires computing $4g$ and checking its projection onto the CF layer.

The code finds: $4g$ projected to $p=3$ (CF layer) gives $[0]$ --- trivial.
But physically, $4g$ should project nontrivially to CF (giving $\mathbb{Z}_8 \to \mathbb{Z}_{16}$).

\subsection{Root cause: the mixed generator has $c_2 = 0$ in $\mathbb{Z}$}

Step-by-step stacking trace reveals the mechanism.
The pip class has $c_2(T,T) = -1$ and the MC class also has $c_2(T,T) = -1$.
Their $\mathbb{Z}$-valued $p=2$ layers are \emph{identical}.
Therefore the mixed generator $g = (-\text{pip}) + \text{MC}$ has $p=2$ layer:
\begin{equation}
g_2(T,T) = -(-1) + (-1) + \underbrace{\mathrm{addTwister}(2,2)}_{=0} = 1 + (-1) = 0.
\end{equation}
The $p=2$ layer of the mixed generator is \textbf{exactly zero as an integer}, not just zero mod 2.

When we stack 4 copies of $g$, the $p=2$ layer stays zero at every step:
\begin{equation}
1g: \; c_2 = 0, \quad 2g: \; c_2 = 0, \quad 3g: \; c_2 = 0, \quad 4g: \; c_2 = 0.
\end{equation}
Since addTwister(3,1) (the CF stacking twister) depends on the $p=2$ layers, and both inputs have $c_2 = 0$, the twister gives zero at $p=3$.
Hence $4g$ has $c_3(T,T,T) = 0$, and the extension $M_{12} \to E_{31}$ is trivial.

In contrast, stacking 2 copies of the pure MC class:
\begin{equation}
1\times\text{MC}: \; c_2(T,T) = -1, \quad 2\times\text{MC}: \; c_2(T,T) = -2.
\end{equation}
At the second step, the addTwister(3,1) sees $c_2$ values $(-1, -1)$ and produces a nontrivial $c_3(T,T,T) = -1$.
After purification, $c_3(T,T,T) = -3$ (odd, nontrivial in $\mathbb{Z}_2$ module).

\subsection{Why is $c_2$ the same for pip and MC?}

The pip basis class has $(n_1, n_2, n_3, \nu_4) = (m_1, 0, 0, 1)$ at the physical $\mathbb{Z}_2$ level.
But its $\mathbb{Z}$-lift representation has $c_2(T,T) = -1$, which is the $\mathbb{Z}$-lift of $1 \in \mathbb{Z}_2$.

The MC basis class has $(n_1, n_2, n_3, \nu_4) = (0, m_1^2, 0, 1)$ at the physical level.
Its $\mathbb{Z}$-lift also has $c_2(T,T) = -1$: the $\mathbb{Z}$-lift of the generator $m_1^2 \in H^2(G_b, \mathbb{Z}_2)$.

Both lifts happen to be the same integer $-1$.
This is \emph{not} a coincidence: the bar resolution for $\mathbb{Z}_2$ has a single generator in degree 2, and both the pip coboundary tail and the MC cocycle map to the same $\mathbb{Z}$-valued basis element.

The consequence: the mixed generator $g = -\text{pip} + \text{MC}$ has its $p=2$ layer cancel in $\mathbb{Z}$, not just in $\mathbb{Z}_2$.
This ``over-cancellation'' kills all subsequent twister contributions.

\subsection{The fix}

The issue is that the $p=2$ layer of the pip class is a ``parasitic'' $\mathbb{Z}$-lift artifact that happens to equal the MC generator's lift.
Physically, the pip class should have $n_2 = 0$ (in $\mathbb{Z}_2$), but the code gives it a $\mathbb{Z}$-lift value $c_2 = -1 \neq 0$ in $\mathbb{Z}$.

There are two possible fixes:
\begin{enumerate}
\item \textbf{Mod out the $\mathbb{Z}$-lift at each layer:}
After constructing the pip class, reduce its $p=2$ layer by the torsion of the coefficient module (i.e., mod 2).
This would set $c_2(T,T)$ to $-1 \bmod 2 = 1$ or to $0$ depending on convention.
However, this breaks the purification mechanism which relies on the $\mathbb{Z}$-lift being consistent across layers.
\item \textbf{Enable addTwister(2,2) and remove the parasitic tail:}
Set the pip class's $p=2$ tail to its ``physical'' value (zero), and let addTwister(2,2) provide the stacking contribution explicitly.
This separates the ``coboundary tail'' from the ``stacking twister'' as in the manual calculation.
\end{enumerate}
The correct fix requires careful thought about the purification algorithm.
This is an open issue.

%%%%%%%%% END TODO: CONTENTS


%%%%%%%%%% TODO: BIBLIOGRAPHY
% Use your bibtex library, formatted by the SciPost style file.
\bibliography{SciPost_Example_BiBTeX_File.bib}

%%%%%%%%%% END TODO: BIBLIOGRAPHY


\end{document}
